%% 1장 기호와 순서 (새 목차, 2026-11-25)
\chapter{Notation and Order}
\label{ch:order}

\abstract{정렬하고, 가장 작은 것을 고르고, 같은 것끼리 묶고, 순위를 매기는 계산은 모두 순서와 동치를 전제한다. 이 장은 동치관계, 약순서, 사전식 순서의 이론에서 출발해 비교 함수가 지켜야 할 계약, 부동소수점 수의 전순서와 정수 키, 안정 정렬과 여러 키의 정렬, 결정적인 동점 규칙, \code{unique}와 순위의 구현까지 따라간다.}


\section{Notation Used in This Book}
\label{sec:notation}

본격적인 이야기에 앞서 책 전체에서 쓸 기호를 한곳에 모아 둔다. 대부분은 해석학과 선형대수학에서 익숙한 것이지만 계산과 맞닿은 몇몇 기호는 분야마다 쓰는 방식이 조금씩 달라서 여기서 한 번 못 박아 두자. 처음 읽을 때는 가볍게 훑어보고 뒤에서 낯선 기호를 만날 때 돌아와 확인해도 충분하다.

\subsection{Numbers and Sets}

자연수의 집합은 $\mathbb{N}=\{1,2,3,\dots\}$로 쓰고, $0$을 포함할 때는 $\mathbb{N}_0=\mathbb{N}\cup\{0\}$으로 쓴다. 정수, 유리수, 실수, 복소수의 집합은 각각 $\mathbb{Z}$, $\mathbb{Q}$, $\mathbb{R}$, $\mathbb{C}$이고, 확장된 실수 $\mathbb{R}\cup\{-\infty,+\infty\}$는 $\overline{\mathbb{R}}$로 쓴다. 컴퓨터가 실제로 다루는 부동소수점 수의 집합은 $\mathbb{F}$로 쓰는데 그 정확한 정의는 부동소수점 수를 다루는 장으로 미룬다.

실수의 구간은 늘 쓰던 대로 $[a,b]$, $[a,b)$, $(a,b]$, $(a,b)$로 쓴다. 한편 컴퓨터 과학에서는 정수의 구간이 자주 필요하므로 정수 $a\le b$에 대해
\[
[a\,..\,b]=\{k\in\mathbb{Z}: a\le k\le b\}
\]
로 쓰기로 한다. 예컨대 길이가 $n$인 배열의 색인 전체는 $[0\,..\,n-1]$이다. 유한집합 $A$의 원소의 개수는 $|A|$로 쓰고, 명제 $P$가 참이면 $1$, 거짓이면 $0$인 값은 Iverson의 괄호를 빌려 $[P]$로 쓴다. 이 괄호를 쓰면 예컨대 $\sum_{i=1}^{n}[x_i<0]$이 음수의 개수가 되어 지시함수를 일일이 정의하지 않아도 세는 식을 간결하게 적을 수 있다.

\subsection{Functions and Operations}

실수 $x$의 바닥과 천장은 $\lfloor x\rfloor=\max\{k\in\mathbb{Z}:k\le x\}$와 $\lceil x\rceil=\min\{k\in\mathbb{Z}:k\ge x\}$이다. 정수의 나눗셈은 프로그래밍 언어마다 약속이 달라서 특히 주의가 필요하다. 이 책에서 $a\bmod b$는 언제나
\[
a\bmod b=a-b\left\lfloor a/b\right\rfloor
\]
를 뜻하며, 따라서 $b>0$이면 $a\bmod b\in[0\,..\,b-1]$이다. Python의 \code{\%} 연산자가 이 약속을 따르는 반면, C와 C++의 \code{\%}는 몫을 $0$ 쪽으로 자르는 약속을 따르므로 $a$와 $b$의 부호가 다르고 $a$가 $b$로 나누어떨어지지 않을 때 두 결과가 다르다. 예컨대 Python에서 \code{-7 \% 3}은 $2$이지만 C에서는 $-1$이다. numpy와 stardust의 정수 나머지는 Python의 약속을 따르므로 C++로 구현할 때는 이 차이를 직접 메워 주어야 한다.

로그는 따로 밑을 적지 않으면 자연로그이고, 밑이 $2$와 $10$인 로그는 각각 $\log_2$와 $\log_{10}$으로 쓴다. 함수 $f$가 집합 $A$ 위에서 가장 작은 값을 갖는 점들의 집합은
\[
\operatorname*{arg\,min}_{x\in A}f(x)=\{x\in A: f(x)\le f(y)\text{ for all }y\in A\}
\]
로 쓴다. 정의에서 보듯 이것은 원소 하나가 아니라 \textbf{집합}이며 원소가 여럿일 수도 있다. 이 집합에서 원소 하나를 골라야 할 때 무엇을 고를지 정하는 일이 이 장 후반의 주제이다.

\subsection{Vectors and Matrices}

벡터는 열벡터로 보며 $x=(x_1,\dots,x_n)^{\mathsf T}\in\mathbb{R}^n$처럼 쓴다. 행렬 $A$의 $(i,j)$ 성분은 $a_{ij}$ 또는 $A_{ij}$로 쓰고, 전치는 $A^{\mathsf T}$, 켤레전치는 $A^{\mathsf H}$로 쓴다. 굵은 글자를 따로 쓰지 않으므로 기호가 스칼라인지 벡터인지는 문맥으로 구별한다. 수학 식에서 색인은 관례대로 $1$부터 세지만 배열과 프로그램을 이야기할 때는 Python과 C++처럼 $0$부터 센다. 두 관례가 한 문단 안에서 섞일 때는 어느 쪽인지 밝혀 둔다.

\subsection{Pseudocode Conventions}

알고리즘은 특정 언어에 매이지 않는 유사코드로 적는다. 대입은 $x\gets y$로 쓰고, \textbf{for} $i\gets a$ \textbf{to} $b$는 $i=a,a+1,\dots,b$에 대해 차례로 도는 반복이다. 배열 $A$의 원소는 $A[i]$로 쓰며 앞에서 말한 대로 색인은 $0$부터 센다. 주석은 \code{//}로 시작한다. 논리 연산 \textbf{and}와 \textbf{or}는 왼쪽부터 평가하고 왼쪽만으로 결과가 정해지면 오른쪽을 평가하지 않는다. 이 약속은 사소해 보이지만 오른쪽 식이 왼쪽 조건이 참일 때만 의미를 갖는 경우(예컨대 $i<n$ \textbf{and} $A[i]=v$)에 유사코드를 짧게 유지해 준다.

\cref{tab:notation}은 이 절의 기호를 한곳에 모은 것이다.

\begin{table}[htbp]
\begin{sdtable}{colspec={lX}, row{1}={font=\small\bfseries}}
\toprule
기호 & 뜻 \\
\midrule
$\mathbb N$, $\mathbb N_0$ & $\{1,2,\dots\}$, $\{0,1,2,\dots\}$ \\
$\overline{\mathbb R}$, $\mathbb F$ & 확장된 실수, 부동소수점 수의 집합(10장) \\
$[a\,..\,b]$ & 정수 구간 $\{k\in\mathbb Z:a\le k\le b\}$ \\
$|A|$, $[P]$ & 원소의 수, Iverson 괄호($P$가 참이면 $1$) \\
$a\bmod b$ & $a-b\lfloor a/b\rfloor$ (Python의 약속) \\
$\operatorname{arg\,min}$ & 최솟값을 갖는 점들의 \textbf{집합} \\
$A^{\mathsf T}$, $A^{\mathsf H}$ & 전치, 켤레전치 \\
$x\gets y$, \code{//} & 유사코드의 대입과 주석; 배열 색인은 $0$부터 \\
\bottomrule
\end{sdtable}
\caption{이 책의 기본 기호.}
\label{tab:notation}
\end{table}

\section{Relations and Orders}
\label{sec:relations}

\subsection{Equivalence Relations and Partitions}

집합 $X$ 위의 \term{이항관계}{binary relation}는 곱집합 $X\times X$의 부분집합 $R$이며, $(a,b)\in R$을 흔히 $a\mathrel{R}b$로 쓴다. 관계가 지닐 수 있는 성질 가운데 이 장에서 쓸 것들을 먼저 모아 두자. $X$의 모든 원소 $a,b,c$에 대해

\begin{itemize}
\item $a\mathrel{R}a$이면 \term{반사적}{reflexive}, 어떤 $a$에 대해서도 $a\mathrel{R}a$가 아니면 \term{비반사적}{irreflexive}이다.
\item $a\mathrel{R}b$이면 $b\mathrel{R}a$일 때 \term{대칭적}{symmetric}이고, $a\mathrel{R}b$와 $b\mathrel{R}a$가 함께 성립하면 $a=b$일 때 \term{반대칭적}{antisymmetric}이다.
\item $a\mathrel{R}b$이고 $b\mathrel{R}c$이면 $a\mathrel{R}c$일 때 \term{추이적}{transitive}이다.
\end{itemize}

\begin{definition}[동치관계]
\label{def:equivalence}
반사적이고 대칭적이며 추이적인 관계를 \term{동치관계}{equivalence relation}라 한다. 동치관계 $\sim$과 $a\in X$에 대해 $[a]=\{b\in X:a\sim b\}$를 $a$의 \term{동치류}{equivalence class}라 하고, 동치류 전체의 집합 $X/{\sim}=\{[a]:a\in X\}$를 \term{몫집합}{quotient set}이라 한다.
\end{definition}

동치관계는 \textbf{같다고 여길 것}을 정하는 일이고 그렇게 같다고 여긴 것끼리 묶으면 집합이 겹치지 않는 조각들로 나뉜다. 그 조각들을 정확히 말하려면 분할이라는 개념이 필요하다. 집합 $X$의 부분집합들의 모임 $\mathcal{P}$가 공집합을 원소로 갖지 않고, 서로 다른 두 원소가 늘 서로소이며, 그 합집합이 $X$이면 $\mathcal{P}$를 $X$의 \term{분할}{partition}이라 한다.

\begin{theorem}
\label{thm:partition}
집합 $X$ 위의 동치관계 $\sim$에 대해 $X/{\sim}$은 $X$의 분할이다. 거꾸로 $X$의 분할 $\mathcal{P}$가 주어지면 ``$a$와 $b$가 $\mathcal{P}$의 같은 원소에 속한다''로 정의한 관계는 동치관계이고, 그 몫집합은 $\mathcal{P}$이다.
\end{theorem}

\begin{proof}
반사성에서 $a\in[a]$이므로 어떤 동치류도 비어 있지 않고, 동치류들의 합집합은 $X$이다. 이제 $[a]\cap[b]\ne\emptyset$이라 하고 $c$를 그 공통 원소로 잡자. 그러면 $a\sim c$이고 $b\sim c$이므로 대칭성과 추이성에서 $a\sim b$이다. 이때 임의의 $x\in[b]$에 대해 $b\sim x$이고, 다시 추이성에서 $a\sim x$이므로 $x\in[a]$, 곧 $[b]\subseteq[a]$이며, 같은 논증을 $a$와 $b$를 바꾸어 적용하면 $[a]=[b]$를 얻는다. 따라서 서로 다른 두 동치류는 서로소이고, $X/{\sim}$은 분할이다.

거꾸로 분할 $\mathcal{P}$가 주어졌다고 하자. 모든 $a$는 $\mathcal{P}$의 원소 가운데 정확히 하나에 속하므로 관계는 반사적이고, 정의가 $a$와 $b$에 대해 대칭이므로 대칭적이다. $a$와 $b$가 $P\in\mathcal{P}$에, $b$와 $c$가 $Q\in\mathcal{P}$에 속하면 $b\in P\cap Q$이므로 서로소 조건에서 $P=Q$이고, 따라서 $a$와 $c$는 같은 원소에 속한다. 마지막으로 $a$의 동치류는 $a$를 포함하는 $\mathcal{P}$의 원소 그 자체이고, $\mathcal{P}$의 원소는 비어 있지 않아 저마다 그 안의 원소의 동치류가 되므로 몫집합은 $\mathcal{P}$이다.
\end{proof}

\cref{thm:partition}는 단순하지만 데이터를 다루는 계산의 상당 부분이 사실상 이 정리의 구현이다. 배열에서 서로 다른 값을 뽑는 \code{unique}, 표의 행을 키로 묶는 \code{groupby}, 값마다 번호를 매기는 \code{factorize}는 모두 ``같은 키''라는 동치관계의 몫집합을 계산하고 각 원소가 어느 동치류에 속하는지를 기록한다. 이런 연산을 올바르게 구현하려면 먼저 ``같다''가 정말 동치관계인지부터 확인해야 한다.

부동소수점 수를 보면 이 확인이 결코 형식적인 절차가 아님을 알 수 있다. 부동소수점 수의 표준인 IEEE~754에서 NaN(not a number)은 자기 자신과도 같지 않다고 비교된다\cite{ieee754-2019}. 곧 \code{==}는 반사적이지 않아 동치관계가 아니므로 묶는 기준으로 그대로 쓰면 \cref{thm:partition}의 전제가 무너져 NaN이 하나도 묶이지 않거나, 구현의 세부에 따라 엉뚱하게 묶인다. 묶는 연산에는 \code{==} 대신 따로 정의한 동치관계를 써야 한다.

\subsection{Partial and Total Orders}

이제 ``같다''가 아니라 ``앞선다''를 다루는 관계로 넘어가자.

\begin{definition}[부분순서와 전순서]
\label{def:partial-order}
반사적이고 반대칭적이며 추이적인 관계 $\le$를 \term{부분순서}{partial order}라 하고, 부분순서가 주어진 집합 $(X,\le)$를 \term{부분순서집합}{partially ordered set, poset}이라 한다. 모든 $a,b\in X$에 대해 $a\le b$이거나 $b\le a$이면, 곧 어떤 두 원소도 비교할 수 있으면 $\le$를 \term{전순서}{total order}라 한다.
\end{definition}

부분순서의 전형적인 예는 셋이다. 첫째로 $\mathbb{N}$ 위의 나누어떨어짐 관계 $a\mid b$는 부분순서이지만 $2$와 $3$처럼 서로 나누지 않는 두 수는 비교할 수 없으므로 전순서가 아니다. 둘째로 어떤 집합의 부분집합들 위의 포함관계 $\subseteq$도 부분순서이다. 셋째로 $\mathbb{R}^n$ 위에서 $x\le y$를 ``모든 $i$에 대해 $x_i\le y_i$''로 정의한 성분별 순서는 부분순서이며, 최적화에서 여러 목표를 함께 비교할 때 쓰는 Pareto 순서가 바로 이것이다. 유한한 부분순서집합은 \cref{fig:hasse}처럼 $a<b$이고 그 사이에 다른 원소가 없을 때만 $a$에서 $b$로 위를 향하는 선을 그은 Hasse 그림으로 나타내면 한눈에 들어온다.

\begin{figure}[htbp]
\centering
\begin{tikzpicture}[every node/.style={font=\small, inner sep=2pt}, x=14mm, y=12mm]
  \node (1) at (0,0) {$1$};
  \node (2) at (-0.8,1) {$2$};
  \node (3) at (0.8,1) {$3$};
  \node (4) at (-1.6,2) {$4$};
  \node (6) at (0,2) {$6$};
  \node (12) at (-0.8,3) {$12$};
  \draw (1)--(2) (1)--(3) (2)--(4) (2)--(6) (3)--(6) (4)--(12) (6)--(12);
\end{tikzpicture}
\caption{$12$의 양의 약수들이 나누어떨어짐 관계로 이루는 부분순서의 Hasse 그림. $4$와 $6$, $2$와 $3$처럼 선을 따라 위로 올라가는 길로 이어지지 않는 두 원소는 비교할 수 없다. $1$과 $12$처럼 길로 이어지면 선이 직접 닿지 않아도 비교할 수 있다.}
\label{fig:hasse}
\end{figure}

부분순서 $\le$가 있으면 $a<b$를 ``$a\le b$이고 $a\ne b$''로 정의할 수 있고, 이렇게 얻은 $<$는 비반사적이고 추이적이다. 거꾸로 비반사적이고 추이적인 관계 $<$가 주어지면 ``$a<b$이거나 $a=b$''로 부분순서를 되찾을 수 있다. 앞의 성질을 갖는 관계를 \term{엄격한 부분순서}{strict partial order}라 하며 이 관계는 자동으로 \textbf{비대칭적}이어서 $a<b$와 $b<a$가 함께 성립하는 일이 없다. 둘이 함께 성립한다면 추이성에서 $a<a$가 되어 비반사성과 모순이 된다. 프로그래밍 언어의 정렬 함수는 거의 예외 없이 ``앞서는가''를 묻는 비교 함수 하나를 받으므로 앞으로 정렬과 비교를 다룰 때는 $\le$보다 $<$를 기본으로 삼는다. 서로 다른 두 원소 $a,b$에 대해 늘 $a<b$이거나 $b<a$인 엄격한 부분순서를 \term{엄격한 전순서}{strict total order}라 한다. 전순서 $\le$에서 얻은 $<$가 그런 예이고, 엄격한 전순서는 동점이 등호뿐인 엄격한 약순서와 같다.

\subsection{Weak Orders}

실제 계산에서 가장 자주 만나는 순서는 전순서가 아니다. 학생들을 키 순으로 세우는 경우를 생각해 보자. 두 학생의 키가 같으면 어느 쪽도 다른 쪽보다 앞서지 않지만 그렇다고 두 학생이 같은 사람인 것은 아니다. ``키가 작다''는 학생들 위의 엄격한 부분순서이되, 비교할 수 없는 두 원소가 \textbf{동점}이라는 특별한 구조를 갖는다. 약순서는 이 구조를 정확히 잡아낸 개념이다.

\begin{definition}[약순서]
\label{def:weak-order}
집합 $X$ 위의 엄격한 부분순서 $<$에 대해 $a\sim b$를 ``$a<b$도 $b<a$도 아니다''로 정의하자. $\sim$이 추이적이면 $<$를 \term{엄격한 약순서}{strict weak ordering}라 하고, $a\sim b$일 때 $a$와 $b$가 \term{동점}{tie}이라고 한다.
\end{definition}

정의의 $\sim$은 언제나 반사적이고 대칭적이므로, 약순서란 비교할 수 없음이 추이적이어서 $\sim$이 동치관계가 되는 엄격한 부분순서라고 말할 수 있다. \cref{fig:hasse}의 나누어떨어짐에서 얻은 엄격한 부분순서(``$a\mid b$이고 $a\ne b$'')는 약순서가 아닌데, $4\sim 3$이고 $3\sim 2$이지만 $2<4$여서 $\sim$이 추이적이지 않다. 반면 학생들의 키는 약순서이며, 다음 정리에 따르면 모든 약순서가 사실상 이런 꼴이다.

\begin{lemma}
\label{lem:weak-mix}
$<$가 집합 $X$ 위의 엄격한 약순서이면 모든 $a,b,c\in X$에 대해 다음이 성립한다.
\begin{enumerate}
\item $a<b$이고 $b\sim c$이면 $a<c$이다.
\item $a\sim b$이고 $b<c$이면 $a<c$이다.
\end{enumerate}
\end{lemma}

\begin{proof}
첫째 주장만 보이면 둘째는 같은 방식으로 따라 나온다. $a<b$이고 $b\sim c$인데 $a<c$가 아니라고 하자. 그러면 $c<a$이거나 $a\sim c$이다. 앞의 경우에는 $c<a<b$에서 추이성으로 $c<b$인데, 이는 $b\sim c$와 모순이다. 뒤의 경우에는 $a\sim c$와 $c\sim b$에서 $\sim$의 추이성으로 $a\sim b$인데, 이는 $a<b$와 모순이다. 따라서 $a<c$이다.
\end{proof}

\begin{theorem}[약순서의 표현]
\label{thm:weak-order-rep}
집합 $X$ 위의 관계 $<$에 대해 다음은 서로 동치이다.
\begin{enumerate}
\item $<$는 엄격한 약순서이다.
\item 어떤 전순서집합 $(T,<_T)$과 함수 $f:X\to T$가 있어 모든 $a,b\in X$에 대해 $a<b$와 $f(a)<_T f(b)$가 동치이다.
\end{enumerate}
이때 $T$로 몫집합 $X/{\sim}$을, $f$로 $a\mapsto[a]$를 잡을 수 있다.
\end{theorem}

\begin{proof}
(ii)에서 (i)로 가자. $<_T$가 비반사적이고 추이적이므로 $<$도 그렇다. 또한 $a\sim b$는 $f(a)<_T f(b)$도 $f(b)<_T f(a)$도 아니라는 뜻이고, $<_T$가 전순서이므로 이는 곧 $f(a)=f(b)$이다. 등호는 추이적이므로 $\sim$도 추이적이다.

(i)에서 (ii)로 가자. $\sim$이 동치관계이므로 몫집합 $T=X/{\sim}$을 만들 수 있다. 동치류 위에서 $[a]<_T[b]$를 $a<b$로 정의하고 싶은데, 그러려면 이 정의가 대표의 선택과 무관해야 한다. $a\sim a'$이고 $b\sim b'$이며 $a<b$라 하자. \cref{lem:weak-mix}의 (ii)를 $a'\sim a<b$에 적용하면 $a'<b$이고, (i)을 $a'<b\sim b'$에 적용하면 $a'<b'$이다. 따라서 $<_T$는 잘 정의된다. $<_T$가 비반사적이고 추이적임은 $<$의 성질에서 바로 나온다. 서로 다른 두 동치류 $[a]\ne[b]$에서는 $a\not\sim b$이므로 $a<b$이거나 $b<a$이고, 따라서 $<_T$는 전순서이다. 마지막으로 $f(a)=[a]$로 두면 정의에서 $a<b$와 $f(a)<_T f(b)$가 동치이다.
\end{proof}

\begin{figure}[htbp]
\centering
\begin{tikzpicture}[x=1cm,y=0.9cm,font=\small]
\foreach \y/\h/\names in {0/{150}/{지아},1/{162}/{민준, 서연},2/{171}/{도윤, 하은, 시우},3/{180}/{예준}}{
  \draw[gray!60] (0,\y) -- (7.2,\y);
  \node[anchor=east] at (-0.1,\y) {$\h$\,cm};
  \node[anchor=west,fill=white] at (0.4,\y) {\names};
}
\draw[->,thick] (8,0) -- (8,3) node[midway,right,align=left] {$<$: 키가\\작은 쪽이 앞선다};
\end{tikzpicture}
\caption{약순서는 키로 묶인 층이다. 같은 줄의 학생들은 서로 동점이고, 줄과 줄 사이는 키의 전순서로 비교된다(\cref{thm:weak-order-rep}).}
\label{fig:weak-layers}
\end{figure}

\cref{thm:weak-order-rep}에 따르면 약순서란 결국 \textbf{키로 비교하는 순서}이다. 원소마다 전순서집합의 값 하나를 키로 붙이고 키가 작은 쪽이 앞선다고 정하면 언제나 약순서가 되며 모든 약순서는 이렇게 얻어진다. 학생과 키, 표의 행과 정렬 기준 열, 점과 기준점까지의 거리가 모두 이런 예이다. 이 정리는 효용 이론에서 선호를 실수값 효용으로 나타내는 문제의 출발점이기도 하다\cite{fishburn1970utility}. 순서 이론 전반에 대해서는 \cite{davey2002order}가 좋은 안내서이다.

약순서는 정렬 알고리즘이 비교 함수에 요구하는 최소한의 조건이기도 하다. C++ 표준은 \code{std::sort}에 넘기는 비교 함수가 엄격한 약순서를 유도해야 한다고 규정하고 이를 어기면 동작을 정의하지 않는다\cite{iso14882-2020}. 부동소수점 수의 \code{<}를 보면 이 조건이 얼마나 쉽게 깨지는지 알 수 있다. IEEE~754에서 NaN과 어떤 수를 \code{<}로 비교하면 언제나 거짓이므로 $1\sim\text{NaN}$이고 $\text{NaN}\sim2$이지만 $1<2$이다. 비교할 수 없음이 추이적이지 않으니 NaN이 섞인 배열에서 \code{<}는 약순서가 아니다.

\subsection{Lexicographic Order}

키가 하나가 아니라 여럿인 경우도 많다. 성적표를 학년순으로, 같은 학년 안에서는 반 순으로, 같은 반 안에서는 이름순으로 정렬하는 것이 그 예이며 이런 비교는 사전에서 낱말을 늘어놓는 방식과 같다.

\begin{definition}[사전식 순서]
\label{def:lex}
집합 $X_1,\dots,X_m$ 위에 엄격한 약순서 $<_1,\dots,<_m$이 있고, 각각의 동점을 $\sim_1,\dots,\sim_m$이라 하자. 곱집합 $X_1\times\cdots\times X_m$의 두 원소 $a=(a_1,\dots,a_m)$와 $b=(b_1,\dots,b_m)$에 대해, 어떤 $k\in[1\,..\,m]$가 있어 $i<k$인 모든 $i$에서 $a_i\sim_i b_i$이고 $a_k<_k b_k$일 때 $a<_{\mathrm{lex}}b$라 하고, $<_{\mathrm{lex}}$를 \term{사전식 순서}{lexicographic order}라 한다.
\end{definition}

정의의 $k$는 $a$와 $b$가 처음으로 동점이 아니게 되는 성분으로, 그 성분에서 앞서는 쪽이 전체로도 앞서고 그 뒤의 성분은 보지 않는다.

\begin{theorem}
\label{thm:lex}
\cref{def:lex}의 사전식 순서 $<_{\mathrm{lex}}$는 엄격한 약순서이고, $a$와 $b$가 동점일 필요충분조건은 모든 $i$에서 $a_i\sim_i b_i$인 것이다. 특히 각 $<_i$가 엄격한 전순서이면 $<_{\mathrm{lex}}$도 엄격한 전순서이다.
\end{theorem}

\begin{proof}
비반사성은 분명하다. $a<_{\mathrm{lex}}a$라면 어떤 $k$에서 $a_k<_k a_k$여야 하기 때문이다.

추이성을 보이자. $a<_{\mathrm{lex}}b$가 성분 $k$에서, $b<_{\mathrm{lex}}c$가 성분 $l$에서 정해졌다고 하고 $p=\min(k,l)$로 두자. $i<p$이면 $a_i\sim_i b_i$이고 $b_i\sim_i c_i$이므로 $\sim_i$의 추이성에서 $a_i\sim_i c_i$이다. 성분 $p$에서는 세 경우를 나누어 본다. $k<l$이면 $a_p<_p b_p$이고 $b_p\sim_p c_p$이므로 \cref{lem:weak-mix}의 (i)에서 $a_p<_p c_p$이고, $k>l$이면 $a_p\sim_p b_p$이고 $b_p<_p c_p$이므로 (ii)에서 $a_p<_p c_p$이며, $k=l$이면 $<_p$의 추이성에서 $a_p<_p c_p$이다. 어느 경우든 $a<_{\mathrm{lex}}c$가 성분 $p$에서 정해진다.

이제 동점의 특징을 보이자. 모든 $i$에서 $a_i\sim_i b_i$이면 어떤 $k$에서도 $a_k<_k b_k$나 $b_k<_k a_k$가 성립하지 않으므로 $a$와 $b$는 동점이다. 거꾸로 어떤 $i$에서 $a_i\not\sim_i b_i$이면 그런 $i$ 가운데 가장 작은 것을 $k$라 할 때 $a_k<_k b_k$이거나 $b_k<_k a_k$이므로, $a$와 $b$는 그 성분에서 비교된다. 따라서 $a$와 $b$가 동점일 필요충분조건은 모든 성분이 동점인 것이고, 성분별 동점이 추이적이므로 사전식 순서의 동점도 추이적이다. 곧 $<_{\mathrm{lex}}$는 엄격한 약순서이다.

마지막으로 모든 $<_i$가 전순서이면 각 $\sim_i$는 등호이므로, 사전식 순서의 동점은 $a=b$뿐이다. 서로 다른 두 원소는 늘 비교되므로 $<_{\mathrm{lex}}$는 전순서이다.
\end{proof}

길이가 서로 다른 유한 수열, 예컨대 문자열에도 사전식 순서를 정의할 수 있다. 두 수열을 앞에서부터 비교하다가 처음으로 다른 자리에서 순서를 정하고 한쪽이 다른 쪽의 앞부분과 완전히 같으면 짧은 쪽을 앞에 두면 된다. 사전에서도 이 규칙에 따라 ``a''가 ``ab''보다, ``ab''가 ``abc''보다 앞에 온다. 이 순서도 전순서이지만 자연수의 순서와 달리 무한히 내려가는 사슬이 있을 수 있어서 조금 낯설다. 실제로
\[
\text{b}>\text{ab}>\text{aab}>\text{aaab}>\cdots
\]
에서 각 문자열은 앞의 것보다 사전식으로 작지만 끝이 없으므로 ``가장 작은 문자열''이 없는 문자열의 집합을 쉽게 만들 수 있다. 길이로 먼저 비교하고 길이가 같을 때만 사전식으로 비교하는 \term{길이 우선 사전식 순서}{shortlex order}에서는 글자의 종류가 유한한 한 이런 일이 생기지 않는다.



\section{Comparators in Practice}
\label{sec:comparators}

\subsection{The Contract of a Comparison Function}

정렬 함수는 원소 두 개를 받아 ``앞의 것이 앞서는가''를 답하는 비교 함수 하나만 안다. 그 함수가 엄격한 약순서를 이룬다는 것이 정렬 알고리즘과 맺는 \term{계약}{contract}이다. 계약을 어긴 비교 함수가 정렬 결과보다 더 많은 것을 망가뜨릴 수 있음은 빠른 정렬 구현의 안쪽 루프를 보면 알 수 있다.

\begin{algorithm}[htbp]
\begin{algorithmic}[1]
\Require 배열 $A[0..n-1]$, 비교 함수 $\mathit{less}$, 기준값 $p$ (구간 어딘가에 있는 원소)
\Ensure $A[0..i-1]$은 $p$보다 앞서고 $A[i]$는 앞서지 않는 첫 원소
\State $i\gets0$
\While{$\mathit{less}(A[i],p)$} \Comment{경계 검사 $i<n$이 없다}
  \State $i\gets i+1$
\EndWhile
\State \Return $i$
\end{algorithmic}
\caption{경계 검사 없는 분할의 안쪽 루프}
\label{alg:unguarded-scan}
\end{algorithm}

\cref{alg:unguarded-scan}은 경계 검사를 하지 않는다. 기준값 $p$가 구간 안에 있고 $\mathit{less}$가 비반사적이면 $\mathit{less}(p,p)$가 거짓이므로 루프는 늦어도 $p$의 자리에서 멈추고 검사 하나를 아낀 만큼 안쪽 루프가 빨라진다. 이 정당화는 약순서의 비반사성 하나에 기댄다. 비교 함수로 \code{<} 대신 \code{<=}를 넘기면 $\mathit{less}(p,p)$가 참이 되어 모든 원소가 같은 배열에서 루프가 구간 끝을 지나 그 너머의 메모리를 읽으며 계속 간다(\cref{fig:unguarded-overrun}). C++ 표준은 엄격한 약순서가 아닌 비교 함수를 넘긴 \code{std::sort}의 동작을 정의하지 않으므로\cite{iso14882-2020} 구현은 이런 최적화를 마음 놓고 쓸 수 있다.

\begin{figure}[htbp]
\centering
\begin{tikzpicture}[x=0.75cm,y=0.75cm,font=\small]
\foreach \k in {0,...,7}{\draw (\k,0) rectangle (\k+1,1); \node at (\k+0.5,0.5) {$5$};}
\foreach \k in {8,9,10}{\draw[dashed,gray] (\k,0) rectangle (\k+1,1); \node[gray] at (\k+0.5,0.5) {?};}
\node[below] at (0.5,0) {$0$}; \node[below] at (7.5,0) {$n-1$};
\draw[->,thick] (0.5,1.6) -- (10.5,1.6);
\node[above] at (5.5,1.6) {$\mathit{less}(5,5)$가 참이면 $i$가 멈추지 않는다};
\draw[decorate,decoration={brace,amplitude=4pt,mirror},yshift=-14pt] (8,0) -- (11,0) node[midway,below=5pt] {배열 밖};
\end{tikzpicture}
\caption{비반사적이지 않은 비교 함수(\code{<=})와 \cref{alg:unguarded-scan}. 원소가 모두 같으면 멈춰야 할 자리가 없어 $i$가 배열 끝을 넘는다.}
\label{fig:unguarded-overrun}
\end{figure}

NaN이 섞인 부동소수점 수의 \code{<}는 비반사성은 지키지만 비교할 수 없음의 추이성을 어긴다(\cref{sec:relations}). 이것도 위험하다. 정렬 구현의 여러 단계가 ``이미 정렬한 구간은 이 비교 함수로 정렬되어 있다''는 귀결을 믿고 경계 검사를 생략하기 때문이다. Kutenin은 2022년에 libc++의 \code{std::sort}를 바꾸면서 실제 코드 가운데 약순서를 어기는 비교 함수가 적지 않아 구현을 바꾸는 일 자체가 그런 코드를 깨뜨린다고 보고했고\cite{kutenin2022sort}, 그 뒤 libc++는 이런 비교 함수가 배열 밖을 읽게 되는 자리마다 검사를 넣었다. stardust의 개발 기록에도 같은 일이 있다. 분위수를 구하는 정렬은 $-0<+0$인 전순서가 필요한데, 조각을 NEON 정렬로 정렬한 뒤 그 전순서로 조각을 병합하게 했더니 쓰기가 배열 밖으로 나갔다. NEON 정렬은 $\pm0$을 같게 보므로 조각 안에서 $-0$과 $+0$이 섞여 있었고, 병합은 ``두 조각이 이 비교 함수로 정렬되어 있다''는 가정 아래 경계 검사 없이 진행했기 때문이다. 고친 방법은 조각마다 영 구간을 $-0$들과 $+0$들로 다시 써서 조각이 정말로 그 전순서로 정렬되게 하는 것이었다. 비교 함수와 이미 정렬된 데이터가 같은 순서를 말해야 한다는 것도 계약의 일부이다.

\subsection{How stardust Makes the Comparison a Weak Order}

stardust의 배열 정렬은 NaN을 비교 함수 안에서 다루지 않는다. 비교할 때마다 NaN인지 묻는 비교 함수는 가지가 늘어 느리고 조건 하나를 빠뜨리면 곧바로 계약을 어긴다. 대신 정렬할 레인을 작업 버퍼로 옮기는 단계에서 NaN을 따로 떼어 낸다.

\begin{algorithm}[htbp]
\begin{algorithmic}[1]
\Require 레인 $x_0,\dots,x_{n-1}$ (띄엄띄엄할 수 있다), 방향
\Ensure 버퍼 $B[0..n-1]$: NaN이 아닌 원소와 NaN이 각각 입력 순서대로 갈라져 있음; NaN의 수 $q$
\State $q\gets\sum_i[x_i\text{가 NaN}]$ \Comment{첫 번째 읽기: 세기만}
\If{오름차순} $p\gets0$; $r\gets n-q$ \Comment{NaN은 가장 큰 값: 끝에}
\Else\ $p\gets q$; $r\gets0$ \Comment{내림차순이면 앞에}
\EndIf
\For{$i=0,\dots,n-1$} \Comment{두 번째 읽기: 한 버퍼에 바로 쓴다}
  \If{$x_i$가 NaN} $B[r]\gets x_i$; $r\gets r+1$
  \Else\ $B[p]\gets x_i$; $p\gets p+1$
  \EndIf
\EndFor
\State \Return $B$, $q$ \Comment{정렬할 것은 NaN이 아닌 $n-q$개뿐}
\end{algorithmic}
\caption{NaN 떼어 내기 (stardust \code{load\_split\_nans})}
\label{alg:split-nans}
\end{algorithm}

\cref{alg:split-nans} 뒤에 남는 원소에는 NaN이 없으므로 평범한 \code{<}가 엄격한 약순서이다(\cref{prop:nan-last}). NaN들은 입력 순서 그대로 한쪽 끝에 놓인다. \cref{prop:nan-last}의 약순서에서 NaN들은 모두 동점이고 안정 정렬은 동점의 입력 순서를 지키므로 이 배치는 안정 정렬의 결과와 같다. 복소수는 \code{unique}에서만 정렬되며 (실수부, 허수부)의 사전식 순서를 쓰고, 실수부나 허수부 가운데 하나라도 NaN이면 그 수 전체를 NaN으로 본다.

\cref{alg:split-nans}은 레인을 두 번 읽는다. 첫 번째 읽기에서 NaN의 수를 알면 두 영역이 어디서 시작하는지 정해지므로, 두 번째 읽기에서 원소마다 제자리에 바로 쓸 수 있고 따로 모았다가 이어 붙이는 복사가 없다. 레인이 띄엄띄엄 놓여 있어도 이 단계가 연속한 작업 버퍼를 만든다. 버퍼는 레인이 짧으면 스레드마다 미리 잡아 둔 64KB 공간을 다시 쓴다. $(100000,8)$ 배열을 축 $1$로 정렬하듯 짧은 레인이 아주 많을 때 레인마다 메모리를 할당하지 않기 위해서이다. \code{argsort}에서는 버퍼의 원소가 값과 원래 색인을 붙인 16바이트 쌍이다.


\section{Ordering Floating-Point Numbers}
\label{sec:float-order}

\subsection{Comparisons in IEEE 754}

부동소수점 수는 실수의 유한한 부분집합이니 실수의 순서를 그대로 물려받으면 될 것 같지만 IEEE~754에는 실수에 없는 값이 셋 더 있다. 무한대 $\pm\infty$, 부호가 다른 두 영 $\pm0$, 그리고 NaN이다. 표준은 부동소수점 수 두 개를 비교한 결과가 ``작다'', ``같다'', ``크다'', ``순서 없음'' 가운데 정확히 하나라고 정하며, 순서 없음은 둘 가운데 적어도 하나가 NaN일 때에만 나온다\cite{ieee754-2019}. 또한 $-0$과 $+0$은 같다고 비교된다. 그래서 \code{<}, \code{<=}, \code{==}는 NaN이 끼면 모두 거짓이고 \code{!=}만 참이다.

이 장의 언어로 옮기면 NaN을 뺀 부동소수점 수의 집합에서 \code{<}는 엄격한 약순서이고 그 동점은 \code{==}이다. 이 동점은 거의 모든 곳에서 등호와 같지만 $-0$과 $+0$이라는 서로 다른 두 비트 패턴을 하나로 묶는다. 반면 앞에서 $1\sim\text{NaN}\sim2$의 예로 확인했듯 NaN을 포함하면 \code{<}는 약순서가 아니므로, NaN이 섞일 수 있는 데이터를 정렬하거나 순위를 매기려면 NaN을 어디에 둘지부터 정해야 한다.

\subsection{Constructing a Total Order}

NaN을 모든 수보다 크다고 보는 것이 가장 흔한 약속으로, numpy와 stardust의 정렬이 이 약속을 따라 오름차순 정렬에서 NaN을 맨 끝에 둔다.

\begin{proposition}
\label{prop:nan-last}
부동소수점 수의 집합 위에서 $a\lessdot b$를 ``$a$와 $b$가 모두 NaN이 아니고 $a<b$이거나, $a$는 NaN이 아니고 $b$는 NaN이다''로 정의하면 $\lessdot$는 엄격한 약순서이다. 그 동치류는 \code{==}로 같은 수들의 묶음과, 모든 NaN을 모은 묶음 하나이다.
\end{proposition}

\begin{proof}
\cref{thm:weak-order-rep}를 쓴다. NaN이 아닌 값 $a$를 그 값이 나타내는 확장된 실수로, 모든 NaN을 $\overline{\mathbb{R}}$에 새로 덧붙인 가장 큰 원소 $\top$로 보내는 함수를 $f$라 하자. $\overline{\mathbb{R}}\cup\{\top\}$은 전순서집합이고, $a\lessdot b$와 $f(a)<f(b)$가 동치임은 정의에서 경우를 나누어 바로 확인된다. 이때 $f$의 값이 같은 원소들이 동치류이므로, $\pm0$은 같은 실수 $0$으로 보내져 한 동치류를 이루고 NaN들은 $\top$으로 보내져 한 동치류를 이룬다.
\end{proof}

\cref{prop:nan-last}의 순서는 여전히 약순서여서 $\pm0$과 여러 NaN 사이의 순서는 정하지 않는다. 비트 패턴까지 구별하는 전순서가 필요하면 IEEE~754가 정의한 \code{totalOrder}를 쓸 수 있다. 이 순서에서 $-0$은 $+0$보다 앞서고 부호가 음인 NaN은 $-\infty$보다 앞서며 부호가 양인 NaN은 $+\infty$보다 뒤에 온다\cite{ieee754-2019}. NaN이 아닌 값만 보면 \code{totalOrder}는 \cref{prop:nan-last}의 순서에서 $-0$을 $+0$보다 앞세운 것과 같다.

\subsection{Mapping to Integer Keys}

정수의 대소는 하드웨어가 한 번에 비교할 수 있고 정수에는 자리별 기수 정렬도 쓸 수 있으므로 각 부동소수점 수를 순서가 같은 부호 없는 정수로 바꾼 뒤 정수로서 정렬하면 부동소수점 수를 빠르게 정렬할 수 있다. 이 사상은 부동소수점 수의 비트 배치 덕분에 가능하다.

binary64 부동소수점 수 $x$의 64비트는 맨 위의 부호 비트 $s\in\{0,1\}$과 나머지 63비트로 이루어지고, 나머지 63비트를 부호 없는 정수로 읽은 값 $m\in[0\,..\,2^{63}-1]$은 NaN이 아닌 한 $|x|$와 같은 순서를 갖는다. 즉 NaN이 아닌 두 수에 대해 $|x|<|y|$와 $m_x<m_y$가 동치이다. 이 성질은 지수를 가수보다 위에 두고 지수를 치우친 정수로 적는 IEEE~754의 설계에서 나온다. 부호를 뺀 63비트는 위의 11비트가 치우친 지수 $e$, 아래 52비트가 가수의 소수부 $f$이고, $e\ge1$이면 값은 $(1+f/2^{52})\,2^{e-1023}$, $e=0$이면 영 또는 비정규 수 $f\cdot2^{-1074}$이다(10장). 지수가 $e$인 수는 모두 $[2^{e-1023},2^{e-1022})$에 있고(비정규 수는 $[0,2^{-1022})$에) 이 구간들은 $e$가 커질수록 위로 겹치지 않고 이어지므로 지수가 큰 수가 크다. 지수가 같으면 값이 $f$에 대해 증가한다. $m=2^{52}e+f$이므로 $m$의 대소는 ($e$, $f$)의 사전식 순서, 곧 값의 크기 순서와 같다. $e=2047$인 무한대와 NaN 가운데 무한대($f=0$)는 모든 유한한 수보다 큰 $m$을 받아 이 순서에 그대로 들어맞는다.

\begin{proposition}
\label{prop:float-key}
NaN이 아닌 부동소수점 수 $x$의 64비트를 부호 없는 정수로 읽은 값을 $u=2^{63}s+m$이라 하고 키를
\[
\kappa(x)=\begin{cases}u+2^{63}=2^{63}+m, & s=0,\\ (2^{64}-1)-u=2^{63}-1-m, & s=1\end{cases}
\]
로 정의하자. 비트로 보면 양수는 부호 비트만 세우고 음수는 모든 비트를 뒤집는다. 그러면 NaN이 아닌 두 수 $x,y$에 대해 \code{totalOrder}에서 $x$가 $y$보다 앞설 필요충분조건은 $\kappa(x)<\kappa(y)$이다. 특히 $\kappa(-0)=2^{63}-1<2^{63}=\kappa(+0)$이고, \code{==}로 같은 두 수는 $\pm0$을 빼면 키도 같다.
\end{proposition}

\begin{proof}
$s=0$인 수의 키는 $[2^{63}\,..\,2^{64}-1]$에, $s=1$인 수의 키는 $[0\,..\,2^{63}-1]$에 놓이므로, 부호가 $1$인 수의 키는 부호가 $0$인 수의 키보다 늘 작다. \code{totalOrder}에서도 부호 비트가 $1$인 값(음수와 $-0$)은 부호 비트가 $0$인 값(양수와 $+0$)보다 늘 앞서므로 두 순서는 부호가 다른 경우에 일치한다.

부호가 같은 두 수를 보자. $s=0$이면 두 수는 $+0$이거나 양수이고 \code{totalOrder}는 크기의 순서와 같으며, 키 $2^{63}+m$은 $m$에 대해 증가하므로 앞에서 보인 성질에 따라 크기의 순서와 같다. $s=1$이면 두 수는 $-0$이거나 음수이고 \code{totalOrder}는 크기의 \textbf{반대} 순서이며, 키 $2^{63}-1-m$은 $m$에 대해 감소하므로 역시 일치한다. 마지막 주장은 $m=0$인 두 영에 정의를 대입한 것이다.
\end{proof}

64비트 정수 $u$의 모든 비트를 뒤집으면 $(2^{64}-1)-u$가 되고 부호 비트를 세우는 것은 $2^{63}$을 더하는 것이므로, $\kappa$는 비트 연산 몇 개로 계산된다. \cref{prop:nan-last}의 약순서로 정렬하려면 NaN을 먼저 따로 떼어 두고 $\pm0$이 같은 키를 갖도록 $-0$을 $+0$으로 바꾼 뒤 $\kappa$를 쓰면 된다. 정수의 비트 표현과 부호 비트를 뒤집는 연산이 왜 이런 성질을 갖는지는 다음 장에서 정수의 나머지 산술과 2의 보수를 다루며 다시 본다.

\cref{tab:float-key}에는 NaN의 키도 있다. NaN까지 포함해도 키가 \code{totalOrder}와 같은 순서를 주는지 묻는 것은 자연스럽다. IEEE~754는 NaN들 사이에서 부호가 음인 NaN을 양인 NaN보다 앞세우고, 양의 NaN에서는 신호 NaN(signaling NaN)을 조용한 NaN(quiet NaN)보다 앞에, 음의 NaN에서는 그 반대로 두며, 나머지(가수 비트의 페이로드) 순서는 구현에 맡긴다\cite{ieee754-2019}.

\begin{proposition}
\label{prop:float-key-nan}
조용한 NaN을 가수의 최상위 비트가 $1$인 것으로 구별하는 IEEE~754의 권장 부호화에서, \cref{prop:float-key}의 $\kappa$를 NaN을 포함한 64비트 패턴 전체에 그대로 쓰면 $\kappa$는 \code{totalOrder}의 순서를 보존한다. 단 같은 부호와 같은 종류(신호 또는 조용한)의 NaN들 사이의 순서는 비트 패턴의 순서로 정한 것이다.
\end{proposition}

\begin{proof}
부호가 다르면 \cref{prop:float-key}의 증명이 그대로 성립한다. 부호가 $0$인 패턴에서 $m$은 유한한 수에서 크기의 순서이고, 지수가 모두 $1$인 패턴 가운데 $+\infty$는 가수가 $0$이라 그다음이며, NaN은 가수가 $0$이 아니어서 $+\infty$보다 큰 $m$을 받는다. NaN끼리는 $m$이 가수의 순서이고 신호 NaN은 가수의 최상위 비트가 $0$이므로 같은 지수의 조용한 NaN보다 작은 $m$을 받는다. 키 $2^{63}+m$이 $m$에 대해 증가하므로 양의 쪽에서 순서가 일치한다. 부호가 $1$인 패턴의 키 $2^{63}-1-m$은 $m$에 대해 감소하므로 $-\infty$가 $-0$보다 작고 음의 NaN이 $-\infty$보다 작으며, 음의 NaN끼리는 조용한 NaN이 신호 NaN보다 작은 키를 받는다. 이것이 \code{totalOrder}가 음의 NaN에 정한 순서이다.
\end{proof}

\subsection{Computing the Key}

\cref{prop:float-key}를 정렬에 쓰려면 NaN과 $\pm0$을 먼저 처리해야 한다. stardust는 두 곳에서 조금 다른 키를 쓴다.

\begin{algorithm}[htbp]
\begin{algorithmic}[1]
\Require NaN이 아닌 binary64 수 $x$, 방향
\Ensure 부호 없는 64비트 키 $k$: 키의 대소가 정렬 순서와 같다
\State \textbf{if} $x=0$ \textbf{then} $x\gets+0$ \Comment{$-0$과 $+0$을 동점으로}
\State $b\gets$ $x$의 비트를 64비트 정수로 읽은 값
\If{$b$의 최상위 비트가 $1$} $u\gets\lnot b$ \Comment{음수: 모든 비트를 뒤집는다}
\Else\ $u\gets b\lor2^{63}$ \Comment{양수: 부호 비트를 세운다}
\EndIf
\State \textbf{if} 내림차순 \textbf{then} $u\gets\lnot u$ \Comment{키를 뒤집으면 순서가 뒤집히고 동점은 그대로}
\State \Return $u$
\end{algorithmic}
\caption{binary64의 정렬 키 (stardust 배열의 \code{radix\_key})}
\label{alg:float-key}
\end{algorithm}

\cref{alg:float-key}의 1행이 없으면 키는 \code{totalOrder}를 따라 $-0$을 $+0$보다 앞세운다. 정렬이 약속하는 것은 \cref{prop:nan-last}의 약순서이므로 두 영을 같은 키로 만들어 동점으로 둔다. 기수 정렬은 키에서 값을 되살리지 않고 원래 값을 키와 함께 옮기므로 안정 정렬이 두 영을 입력 순서대로, 부호를 지닌 채 남긴다. 키에서 값을 되살린다면 $-0$이 모두 $+0$이 된다. 6행은 내림차순을 키 하나로 처리한다. 64비트 정수 $u$에서 $\lnot u=(2^{64}-1)-u$는 감소 함수이므로 순서가 정확히 뒤집히고 같은 키는 같은 키로 가므로 동점은 동점으로 남는다. 결과를 뒤집는 것과 달리 안정성이 깨지지 않는다.

\cref{tab:float-key}는 특수값들의 비트와 키이다(\code{figures/01/keymap.py}로 확인했다). 키는 \code{totalOrder}의 순서대로 엄격히 늘어난다. 음의 NaN이 가장 작은 키를, 양의 NaN이 가장 큰 키를 받는 것도 \code{totalOrder}와 같다. stardust는 NaN을 미리 떼어 내므로 이 두 끝을 정렬에 쓰지 않지만 데이터프레임의 키는 다르게 처리한다. 데이터프레임의 \code{sort\_values}와 \code{groupby}의 그룹 순서 정렬은 모든 NaN을 하나의 NaN으로 바꾼 뒤 같은 사상을 써서 NaN이 $+\infty$ 위의 한 키에 모이게 한다. 결측값(NA)은 키가 아니라 따로 떼어 두었다가 키마다 정한 쪽 끝에 놓는다.

\begin{table}[htbp]
\begin{sdtable}{colspec={lll}, row{1}={font=\small\bfseries}}
\toprule
값 & 비트 (16진) & 키 (16진) \\
\midrule
$-\mathrm{NaN}$ & \code{FFF8000000000000} & \code{0007FFFFFFFFFFFF} \\
$-\infty$ & \code{FFF0000000000000} & \code{000FFFFFFFFFFFFF} \\
$-1$ & \code{BFF0000000000000} & \code{400FFFFFFFFFFFFF} \\
가장 작은 음의 비정규 수 & \code{8000000000000001} & \code{7FFFFFFFFFFFFFFE} \\
$-0$ & \code{8000000000000000} & \code{7FFFFFFFFFFFFFFF} \\
$+0$ & \code{0000000000000000} & \code{8000000000000000} \\
가장 작은 양의 비정규 수 & \code{0000000000000001} & \code{8000000000000001} \\
$1$ & \code{3FF0000000000000} & \code{BFF0000000000000} \\
$+\infty$ & \code{7FF0000000000000} & \code{FFF0000000000000} \\
$+\mathrm{NaN}$ & \code{7FF8000000000000} & \code{FFF8000000000000} \\
\bottomrule
\end{sdtable}
\caption{binary64 특수값의 비트와 \cref{alg:float-key}의 키(1행의 영 처리 전). 키는 위에서 아래로 엄격히 늘어난다.}
\label{tab:float-key}
\end{table}

\begin{figure}[htbp]
\centering
\begin{tikzpicture}[x=1cm,y=1cm,font=\small]
\draw[->] (0,0) -- (10.6,0) node[right] {키};
\foreach \x/\l in {0.3/{$-\mathrm{NaN}$},1.2/{$-\infty$},2.6/{$-1$},4.3/{$-0$},5.3/{$+0=2^{63}$},7.0/{$1$},8.6/{$+\infty$},9.6/{$+\mathrm{NaN}$}}{\draw (\x,0.1) -- (\x,-0.1); \node[above] at (\x,0.12) {\l};}
\draw[decorate,decoration={brace,amplitude=4pt,mirror}] (0.1,-0.35) -- (4.6,-0.35) node[midway,below=5pt] {부호 비트 $1$: 모든 비트를 뒤집음};
\draw[decorate,decoration={brace,amplitude=4pt,mirror}] (5.0,-0.35) -- (9.9,-0.35) node[midway,below=5pt] {부호 비트 $0$: 부호 비트를 세움};
\end{tikzpicture}
\caption{키의 수직선. 음수는 크기가 클수록 작은 키를 받아 $[0,2^{63})$에, 양수는 크기가 클수록 큰 키를 받아 $[2^{63},2^{64})$에 놓인다.}
\label{fig:float-key-line}
\end{figure}


\section{Sorting and Ties}
\label{sec:sorting}

\subsection{Definition of Sorting}

정렬은 누구나 아는 개념처럼 보이지만 동점이 있으면 무엇이 올바른 답인지를 먼저 정해야 한다. 수열 $x_1,\dots,x_n\in X$과 $X$ 위의 엄격한 약순서 $<$가 주어졌다고 하자.

\begin{definition}[정렬]
\label{def:sorting}
$[1\,..\,n]$의 순열 $\pi$가 $i<j$인 모든 $i,j$에 대해 $x_{\pi(j)}<x_{\pi(i)}$를 만족하지 않을 때 $\pi$를 \term{정렬 순열}{sorting permutation}이라 하고, 수열 $x_{\pi(1)},\dots,x_{\pi(n)}$을 정렬된 수열이라 한다.
\end{definition}

정의는 ``뒤에 오는 원소가 앞의 원소보다 엄격히 작은 일이 없다''는 것만 요구할 뿐 동점인 원소의 순서에 대해서는 아무것도 말하지 않는다. 정렬 순열은 언제나 존재한다. $\{x_1,\dots,x_n\}$에서 다른 어떤 원소보다도 엄격히 크지 않은 원소인 극소 원소를 하나 골라 맨 앞에 두고 나머지에 대해 같은 일을 되풀이하면 된다. 극소 원소가 없다면 $x_{i_1}>x_{i_2}>\cdots$인 끝없는 사슬이 유한집합 안에 생겨 어떤 원소가 두 번 나타나고 추이성에서 그 원소가 자기 자신보다 작아지므로, 극소 원소는 늘 있다. 이렇게 만든 순서에서 뒤의 원소가 앞의 원소보다 엄격히 작다면 앞의 원소를 고를 때 그것이 극소가 아니었다는 뜻이니 모순이다.

정렬 순열은 일반적으로 하나가 아니지만 동점을 같은 것으로 보면 정렬된 수열은 하나뿐이다.

\begin{proposition}
\label{prop:sorted-unique}
$\pi$와 $\sigma$가 모두 정렬 순열이면 모든 $i$에 대해 $x_{\pi(i)}\sim x_{\sigma(i)}$이다. 더 구체적으로, 동치류 $C$에 속하는 원소들은 어느 정렬 순열에서나 연속한 위치 $c+1,\dots,c+e$를 차지한다. 여기서 $c$는 $C$의 원소보다 엄격히 작은 $x_k$의 개수이고 $e$는 $C$에 속하는 $x_k$의 개수이다.
\end{proposition}

\begin{proof}
정렬 순열 $\pi$와 $C$의 원소 $x_{\pi(i)}$를 하나 잡자. $x_k<x_{\pi(i)}$인 원소는 정렬의 정의에서 $x_{\pi(i)}$보다 뒤에 올 수 없으므로 모두 앞에 있고, $x_{\pi(i)}<x_k$인 원소는 같은 이유로 모두 뒤에 있다. \cref{lem:weak-mix}에서 ``$C$의 원소보다 엄격히 작다''는 대표의 선택과 무관하므로, $C$의 모든 원소는 엄격히 작은 원소 $c$개보다 뒤에, 엄격히 큰 원소들보다 앞에 놓인다. 남은 자리는 $c+1,\dots,c+e$뿐이고 $C$의 원소가 정확히 $e$개이므로 이 자리를 모두 차지한다. 이 서술은 $\pi$에 의존하지 않으므로 $\sigma$에서도 같고, 따라서 같은 위치에는 같은 동치류의 원소가 온다.
\end{proof}

\subsection{Stable Sorting}

동점인 원소들 사이의 순서는 원래 순서를 그대로 지켜 정하는 것이 가장 자연스럽다.

\begin{definition}[안정 정렬]
\label{def:stable}
정렬 순열 $\pi$가 $i<j$이고 $x_{\pi(i)}\sim x_{\pi(j)}$일 때마다 $\pi(i)<\pi(j)$를 만족하면 $\pi$를 \term{안정 정렬 순열}{stable sorting permutation}이라 한다. 결과가 늘 안정 정렬 순열인 정렬 알고리즘을 \term{안정 정렬}{stable sort}이라 한다.
\end{definition}

안정 정렬의 결과는 동점이 있더라도 하나로 정해진다. 원소의 원래 위치를 둘째 키로 붙여 전순서를 만들어 보면 이 사실이 가장 깔끔하게 드러난다.

\begin{theorem}
\label{thm:stable}
색인 $[1\,..\,n]$ 위의 관계 $\prec$를
\[
i\prec j\iff x_i<x_j\text{ 이거나, }x_i\sim x_j\text{ 이고 }i<j
\]
로 정의하면 $\prec$는 엄격한 전순서이다. 또한 순열 $\pi$가 $<$에 대한 안정 정렬 순열일 필요충분조건은 $\pi(1)\prec\pi(2)\prec\cdots\prec\pi(n)$인 것이다. 따라서 안정 정렬 순열은 정확히 하나 존재한다.
\end{theorem}

\begin{proof}
$\prec$는 쌍 $(x_i,i)$의 사전식 순서이다. 첫 성분의 순서 $<$는 엄격한 약순서이고 둘째 성분의 순서는 정수 위의 엄격한 전순서이므로, \cref{thm:lex}에서 $\prec$는 엄격한 약순서이고 $i$와 $j$가 동점일 필요충분조건은 $x_i\sim x_j$이고 $i=j$인 것이다. 즉 서로 다른 두 색인은 늘 비교되므로 $\prec$는 엄격한 전순서이다.

$\pi$가 안정 정렬 순열이라 하고 $i<j$를 잡자. 정렬의 정의에서 $x_{\pi(j)}<x_{\pi(i)}$가 아니므로 $x_{\pi(i)}<x_{\pi(j)}$이거나 둘이 동점이다. 앞의 경우는 $\pi(i)\prec\pi(j)$를 뜻하고, 뒤의 경우에는 안정성에서 $\pi(i)<\pi(j)$이므로 역시 $\pi(i)\prec\pi(j)$이다. 거꾸로 $\pi(1)\prec\cdots\prec\pi(n)$이라 하고 $i<j$를 잡으면 $\pi(i)\prec\pi(j)$인데, 정의에서 이는 $x_{\pi(j)}<x_{\pi(i)}$를 배제하고, 둘이 동점이면 $\pi(i)<\pi(j)$를 뜻한다. 곧 $\pi$는 안정 정렬 순열이다. 유한한 전순서집합의 원소를 오름차순으로 늘어놓는 방법은 하나뿐이므로 안정 정렬 순열도 하나뿐이다.
\end{proof}

\cref{thm:stable}에서 실용적인 사실 두 가지를 얻는다. 첫째, 안정 정렬은 \textbf{불안정 정렬에 색인을 둘째 키로 붙인 것}과 같으므로 안정성을 보장하지 않는 빠른 알고리즘으로도 (값, 색인) 쌍을 정렬하면 안정 정렬의 결과를 얻는다. 둘째, 동점이 남아 있는 한 불안정 정렬의 결과는 구현과 입력의 세부, 그리고 병렬 실행이라면 스레드의 수와 일의 나뉨에 따라 달라질 수 있으므로 정렬의 답을 하나로 못 박고 싶다면 언제나 전순서를 쓰면 된다. Python의 내장 정렬은 안정성을 보장하며\cite{python-sorting-howto}, 많은 라이브러리도 같은 이유로 기본 정렬을 안정 정렬로 둔다.

\cref{prop:sorted-unique}에서 보듯 정렬된 결과로 \textbf{값만} 내놓고 동점이 곧 같은 값을 뜻한다면, 같은 값은 어느 순서로 놓아도 같은 값이므로 안정 여부는 결과에 드러나지 않는다. 그런데 부동소수점 수에는 동점이지만 구별되는 값이 있다. $-0.0$과 $+0.0$은 \code{<}로 비교하면 동점이지만 비트가 다르고, $1/(-0.0)=-\infty$와 $1/(+0.0)=+\infty$처럼 계산 결과를 바꿀 수도 있다.

\subsection{Stability in stardust's Array Sort}

stardust의 \code{sort}와 \code{argsort}는 기본이 안정 정렬이다(\code{stable=True}). 정렬 하나를 모든 길이에 쓰지 않고 레인의 길이와 dtype으로 알고리즘을 고르는데, 고르는 기준은 데이터를 보지 않고 길이만 본다. \cref{tab:sort-paths}은 그 경로들과 각 경로에서 안정성과 $\pm0$이 지켜지는 이유이다. 경계값들은 측정으로 정했고 알고리즘 자체와 측정은 18장과 31장에서 다룬다.

\begin{table}[htbp]
\begin{sdtable}{colspec={lllX}, row{1}={font=\small\bfseries}}
\toprule
대상 & 길이 & 알고리즘 & 안정성과 $\pm0$ \\
\midrule
값, 쌍 & $\le16$ & Batcher 정렬 네트워크 & 교환할 때 값이 같으면 원래 색인을 비교한다: 안정 \\
값 & $17$--$31$ & 합병 정렬 & 같으면 왼쪽을 먼저 가져온다: 안정 \\
값 & $32$--$1023$ & NEON 정렬 & 불안정; \code{float64}는 영의 부호를 되돌린다(\cref{alg:zero-signs}) \\
값 & $\ge1024$ & LSD 기수 정렬 & 계수 정렬의 반복: 안정; $-0$은 $+0$의 키 \\
쌍 & $17$--$127$ & 합병 정렬(\code{stable=False}면 introsort) & 합병은 안정, introsort는 불안정 \\
쌍 & $\ge128$ & LSD 기수 정렬 & 안정(\code{stable}과 상관없이) \\
\bottomrule
\end{sdtable}
\caption{stardust 배열 정렬의 경로(한 스레드 기준; 긴 레인의 병렬은 \cref{sec:long-lane}). ``쌍''은 \code{argsort}와 \code{unique}가 정렬하는 (값, 원래 색인)의 쌍이다. \code{int64}의 키는 부호 비트만 뒤집은 것이다.}
\label{tab:sort-paths}
\end{table}

쌍을 정렬할 때 비교 함수는 값만 본다. 그러므로 쌍의 정렬이 안정한 것은 \cref{thm:stable}의 $\prec$로 비교해서가 아니라 고른 알고리즘이 안정하기 때문이다. 정렬 네트워크만 예외적으로 값이 같으면 원래 색인을 비교하는데, 정렬 네트워크는 비교와 교환의 위치가 입력과 상관없이 정해져 있어서 동점의 순서를 지킬 방법이 그것뿐이기 때문이다. 이 구분은 \code{stable=False}의 뜻을 정한다. \code{stable=False}는 ``정렬된 순열이기만 하면 된다''는 뜻이고, 쌍의 길이가 $17$--$127$이면 introsort가 동점의 순서를 흩뜨릴 수 있다. 반대로 \code{float64}의 값 정렬은 이 플래그를 무시하고 늘 안정한 경로를 탄다. 값만 정렬할 때는 안정한 결과도 불안정한 결과의 하나이고, 안정한 경로가 더 빠르기 때문이다.

계수 정렬이 안정하다는 것은 한 줄로 보일 수 있다. 원소를 입력 순서대로 읽어 자기 키의 통에서 다음 빈 칸에 넣으므로, 키가 같은 두 원소는 먼저 읽힌 쪽이 먼저 자리를 받는다. 기수 정렬은 바이트마다 계수 정렬을 하는 것이고 \cref{thm:lsd}로 전체가 안정하다.

값만 정렬하는 \code{sort}에서 길이 $32$--$1023$의 NEON 정렬은 안정하지 않다. 그래도 결과는 안정 정렬과 같아야 한다. \cref{prop:sorted-unique}가 길을 알려 준다. 값만 내놓는 정렬에서 동점인 원소들은 연속한 구간을 차지하고, 같은 값이면 어느 순서로 놓아도 결과가 같다. NaN을 미리 떼어 낸 \code{float64}에서 동점이지만 구별되는 값은 $-0$과 $+0$뿐이다. 그러므로 두 영의 부호만 입력 순서대로 되돌려 놓으면 된다.

\begin{algorithm}[htbp]
\begin{algorithmic}[1]
\Require NaN이 아닌 \code{float64} 버퍼 $a_0,\dots,a_{m-1}$, 방향
\Ensure $a$가 안정 정렬의 결과와 비트까지 같게 정렬됨
\State $z\gets0$; $\mathit{neg}\gets0$; 비트열 $w\gets0$
\For{$i=0,\dots,m-1$} \Comment{가지 없이: 영이 아니면 아무것도 쓰지 않는 셈}
  \State $e\gets[a_i=0]$; $s\gets$ $a_i$의 부호 비트
  \State $w_z\gets w_z\lor(s\land e)$; $z\gets z+e$; $\mathit{neg}\gets\mathit{neg}+(s\land e)$
\EndFor
\State NEON 정렬로 $a$를 오름차순 정렬; 내림차순이면 뒤집는다
\If{$0<\mathit{neg}<z$} \Comment{두 부호가 섞였을 때만}
  \State $t\gets$ 영 구간의 시작 \Comment{오름차순: 첫 $a_t\ge0$, 내림차순: 첫 $a_t\le0$}
  \For{$k=0,\dots,z-1$}
    \State $a_{t+k}\gets(w_k=1)\ ?\ -0:+0$
  \EndFor
\EndIf
\end{algorithmic}
\caption{안정하지 않은 정렬 뒤 영의 부호 되돌리기 (stardust \code{zero\_signs})}
\label{alg:zero-signs}
\end{algorithm}

\cref{alg:zero-signs}가 안정 정렬과 같은 결과를 낸다는 것은 \cref{prop:sorted-unique}에서 바로 나온다. 정렬이 끝나면 영들은 어떤 정렬로든 같은 자리 $t,\dots,t+z-1$을 차지하고, 안정 정렬은 그 자리에 영들을 입력 순서대로 놓는다. 부호 말고는 구별할 것이 없으므로 입력 순서대로 기록한 부호를 그 자리에 쓰면 그것이 안정 정렬의 결과이다. 같은 이유로 내림차순에서는 결과를 뒤집어도 된다. 동점인 원소들의 상대 순서는 뒤집히지만 영의 부호를 다시 쓰므로 드러나지 않는다. 부호 기록은 원소마다 비교와 비트 연산 몇 개를 더하지만 가지가 없어 정렬 네트워크의 시간에 견주면 작고 부호가 한쪽뿐이면(흔한 경우) 되쓰기 자체를 건너뛴다. 비교 함수를 쓰는 경로(\code{argsort}의 쌍 정렬, 합병 정렬)에서는 내림차순을 만들 때 결과를 뒤집지 않는다. 대신 값의 비교만 뒤집고, 안정한 알고리즘이 값이 같은 원소들의 입력 순서를 지키게 둔다. 결과를 뒤집으면 동점의 색인이 내림차순으로 놓이기 때문이다.

\begin{figure}[htbp]
\centering
\begin{tikzpicture}[x=0.8cm,y=0.8cm,font=\small]
\node[anchor=east] at (-0.3,2.5) {입력};
\foreach \v [count=\i from 0] in {3,{-0},{-2},{+0},1,{-0}}{\draw (\i,2) rectangle (\i+1,3); \node at (\i+0.5,2.5) {$\v$};}
\node[anchor=west] at (6.3,2.5) {영의 부호: $1,0,1$};
\node[anchor=east] at (-0.3,1) {NEON 정렬};
\foreach \v [count=\i from 0] in {{-2},{+0},{-0},{-0},1,3}{\draw (\i,0.5) rectangle (\i+1,1.5); \node at (\i+0.5,1) {$\v$};}
\fill[gray!25] (1,-1) rectangle (4,0);
\node[anchor=east] at (-0.3,-0.5) {되쓰기 뒤};
\foreach \v [count=\i from 0] in {{-2},{-0},{+0},{-0},1,3}{\draw (\i,-1) rectangle (\i+1,0); \node at (\i+0.5,-0.5) {$\v$};}
\end{tikzpicture}
\caption{\cref{alg:zero-signs}의 예. NEON 정렬은 세 영의 순서를 바꿔 놓지만, 입력 순서대로 기록한 부호 $1,0,1$을 영 구간(회색)에 다시 쓰면 안정 정렬의 결과 $-0,+0,-0$이 된다.}
\label{fig:zero-signs}
\end{figure}


\subsection{Sorting a Long Lane in Parallel}
\label{sec:long-lane}

길이가 $2^{17}$ 이상인 레인 하나는 스레드 풀이 있을 때 여러 스레드로 정렬한다. 나누는 방법은 길이로만 정한다. 레인을 $2^{16}$개씩의 조각으로 자르고 조각마다 \cref{tab:sort-paths}의 직렬 정렬을 한 뒤, 정렬된 조각들을 두 개씩 병합하기를 층마다 되풀이한다. 병합 하나도 여러 스레드로 나눈다. 출력을 $2^{15}$개씩의 덩어리로 자르고, 덩어리가 두 런의 어디서 시작하는지를 이분 탐색으로 찾는다.

\begin{algorithm}[htbp]
\begin{algorithmic}[1]
\Function{CoRank}{$k$, 런 $A[0..n_A)$, 런 $B[0..n_B)$}
  \State $\mathit{lo}\gets\max(0,k-n_B)$; $\mathit{hi}\gets\min(k,n_A)$
  \While{$\mathit{lo}<\mathit{hi}$}
    \State $i\gets\lfloor(\mathit{lo}+\mathit{hi})/2\rfloor$; $j\gets k-i$
    \If{$j>0$ \textbf{and} $i<n_A$ \textbf{and not} $\mathit{less}(B[j-1],A[i])$}
      \State $\mathit{lo}\gets i+1$ \Comment{$A[i]$가 $B[j-1]$보다 먼저 나간다: $A$를 더 가져온다}
    \Else\ $\mathit{hi}\gets i$
    \EndIf
  \EndWhile
  \State \Return $\mathit{lo}$ \Comment{병합 결과의 처음 $k$개 가운데 $A$에서 온 것의 수}
\EndFunction
\ForAll{출력 덩어리 $[k_0,k_1)$ \textbf{in parallel}}
  \State $i_0\gets$\Call{CoRank}{$k_0$}; $i_1\gets$\Call{CoRank}{$k_1$}; $j_0\gets k_0-i_0$; $j_1\gets k_1-i_1$
  \State $A[i_0..i_1)$와 $B[j_0..j_1)$를 병합해 출력 $[k_0,k_1)$에 쓴다 \Comment{같으면 $A$ 먼저}
\EndFor
\end{algorithmic}
\caption{덩어리로 나눈 병합(merge path)}
\label{alg:merge-path}
\end{algorithm}

이 방법이 직렬 안정 정렬과 같은 결과를 내는 것은 두 단계로 보인다. 첫째, 안정하게 정렬된 두 런 $A$(입력의 앞부분)와 $B$(그 뒤)를 ``값이 같으면 $A$ 먼저''로 병합하면 결과는 두 런을 합친 구간의 안정 정렬이다. 병합 결과에서 두 원소의 순서를 보면, 값이 다르면 값의 순서이고, 값이 같고 같은 런에서 왔으면 그 런의 순서(곧 입력 순서)이며, 값이 같고 다른 런에서 왔으면 $A$의 원소, 곧 입력에서 앞선 원소가 먼저이다. 이것이 \cref{thm:stable}의 $\prec$의 순서이다. 둘째, \code{CoRank}$(k)$는 이 병합 결과의 처음 $k$개 가운데 $A$에서 온 원소의 수 $i$를 이분 탐색으로 찾는다. 그런 $i$에서는 $A[i-1]$이 $B[k-i]$보다 앞서고 $B[k-i-1]$이 $A[i]$보다 앞서야 하며, 루프의 조건이 바로 이 두 부등식을 시험한다. 덩어리마다 시작과 끝을 이렇게 정하면 덩어리들은 직렬 병합의 출력을 겹치지 않게 나눠 쓴다. 그러므로 결과는 스레드 수와 상관없이 직렬 병합과 같다.

길이가 $2^{17}$ 이상인 \code{int64}와 \code{float64}의 기수 정렬은 조각 병합 대신 통과 하나하나를 병렬로 한다. 조각마다 이번 바이트의 히스토그램을 세고, (바이트 값, 조각)의 순서로 누적해 조각마다 바이트 값별 첫 자리를 정한 뒤, 조각마다 자기 원소를 그 자리부터 흩뿌린다. 원소 하나가 받는 자리는 ``자기보다 작은 바이트 값의 원소 수 전체 + 같은 바이트 값을 가진 앞 조각들의 원소 수 + 자기 조각 안에서 같은 바이트 값을 가진 앞 원소의 수''이고, 직렬 계수 정렬이 조각을 차례로 읽으며 주는 자리와 정확히 같다.

불안정한 정렬만이 조각을 결과에 드러낸다. introsort가 동점의 순서를 어떻게 흩뜨리는지는 조각의 경계에 달려 있기 때문이다. 그래서 \code{stable=False}인 긴 레인은 풀이 없어도 늘 조각으로 나누어 정렬하고, 결과는 조각 경계라는 길이의 함수로만 정해져 스레드 수와 무관하다. ``나누는 방법을 길이로만 정한다''는 원칙이 여기서 결정성을 지킨다.

\subsection{Sorting by Multiple Keys}

\cref{thm:lex}와 \cref{thm:stable}를 합치면 여러 키로 정렬하는 오래된 기법 하나를 증명할 수 있다. 키 $m$개를 사전식으로 비교해 정렬하고 싶을 때 가장 덜 중요한 키부터 가장 중요한 키까지 차례로 \textbf{안정 정렬을 되풀이}하면 된다는 것이다(\cref{alg:multikey}).

\begin{algorithm}[htbp]
\begin{algorithmic}[1]
\Require 수열 $x_0,\dots,x_{n-1}$과 키 함수 $k^{(1)},\dots,k^{(m)}$, 각 키 위의 엄격한 약순서
\Ensure 키 $(k^{(1)},\dots,k^{(m)})$의 사전식 순서에 대한 안정 정렬 순열 $\pi$
\State $\pi\gets(0,1,\dots,n-1)$ \Comment{항등 순열에서 시작한다}
\For{$j\gets m$ \textbf{downto} $1$}
  \State $\pi\gets$ $\pi$의 원소들을 키 $k^{(j)}(x_{\pi(\cdot)})$로 안정 정렬한 순서
\EndFor
\State \Return $\pi$
\end{algorithmic}
\caption{덜 중요한 키부터 되풀이하는 안정 정렬}
\label{alg:multikey}
\end{algorithm}

\begin{figure}[htbp]
\centering
\begin{tikzpicture}[x=1.2cm,y=0.55cm,font=\small]
\def\col#1#2#3{\node[font=\small\bfseries] at (#1,1.3) {#2};
  \foreach \r [count=\i from 0] in {#3}{\node at (#1,-\i) {\r};}}
\col{0}{입력}{{(2, 나)},{(1, 다)},{(2, 가)},{(1, 가)},{(2, 다)}}
\col{2.6}{둘째 키로}{{(2, 가)},{(1, 가)},{(2, 나)},{(1, 다)},{(2, 다)}}
\col{5.2}{첫째 키로}{{(1, 가)},{(1, 다)},{(2, 가)},{(2, 나)},{(2, 다)}}
\draw[->] (0.85,-2) -- (1.75,-2); \draw[->] (3.45,-2) -- (4.35,-2);
\end{tikzpicture}
\caption{두 키 (학년, 반)의 사전식 순서를 \cref{alg:multikey}로 정렬하는 과정. 덜 중요한 둘째 키로 먼저 안정 정렬하고 첫째 키로 다시 안정 정렬하면, 첫째 키가 같은 행들은 앞 단계의 둘째 키 순서를 지킨다.}
\label{fig:lsd-steps}
\end{figure}

\begin{theorem}
\label{thm:lsd}
\cref{alg:multikey}가 돌려주는 $\pi$는 사전식 순서 $(k^{(1)},\dots,k^{(m)})$에 대한 안정 정렬 순열이다.
\end{theorem}

\begin{proof}
$j=m,m-1,\dots,1$에 대해, $j$번째 단계가 끝난 뒤의 $\pi$가 키 $(k^{(j)},\dots,k^{(m)})$의 사전식 순서에 대한 안정 정렬 순열임을 귀납법으로 보인다. \cref{thm:stable}에 따라 이는 색인 위의 전순서 $\prec_j$의 오름차순으로 $\pi$가 늘어서 있다는 뜻이다. 여기서 $a\prec_j b$는 키 벡터 $(k^{(j)}(x_a),\dots,k^{(m)}(x_a))$가 $(k^{(j)}(x_b),\dots,k^{(m)}(x_b))$보다 사전식으로 앞서거나, 두 키 벡터가 동점이고 $a<b$라는 뜻이다. $j=m$이면 이는 \cref{alg:multikey}의 첫 단계가 항등 순열을 $k^{(m)}$으로 안정 정렬한다는 사실 그 자체이다.

$j+1$에서 성립한다고 하고, $j$번째 단계의 결과에서 두 색인 $a,b$의 상대 순서를 보자. $k^{(j)}(x_a)$와 $k^{(j)}(x_b)$가 동점이 아니면, 이번 정렬은 둘을 $k^{(j)}$의 순서대로 놓는데, 이는 사전식 순서가 첫 성분 $k^{(j)}$에서 정해진다는 것과 일치한다. 둘이 동점이면 이번 정렬이 안정적이므로 둘의 순서는 직전 단계의 순서, 곧 귀납 가정에서 $\prec_{j+1}$의 순서와 같다. 그런데 첫 성분이 동점일 때 $\prec_j$는 정확히 $\prec_{j+1}$과 같다. 따라서 어느 경우든 두 색인은 $\prec_j$의 순서대로 놓이고, 귀납이 완성된다. $j=1$일 때가 정리의 주장이다.
\end{proof}

이 정리에서 각 단계는 \textbf{안정적}이어야 하며 한 단계라도 불안정하면 그 단계에서 동점인 원소들이 앞 단계에서 얻은 순서를 잃어버린다. 또 가장 중요한 키를 \textbf{마지막에} 정렬하므로 순서가 직관과 반대이다. 이 원리는 수를 한 자리씩 정렬하는 \term{기수 정렬}{radix sort}의 바탕이 된다. 자릿수를 맞추어 앞을 $0$으로 채운 음이 아닌 정수에서는 각 자리를 키로 보면 수의 대소가 높은 자리부터 비교하는 사전식 순서이므로 낮은 자리부터 한 자리씩 안정 정렬하면 전체가 정렬된다. 기수 정렬과 그 비용은 정렬을 다루는 장에서 자세히 본다.

\begin{aside}[천공 카드 정렬기]
\cref{thm:lsd}의 방법은 컴퓨터보다 먼저, 천공 카드를 다루던 기계식 정렬기에서 쓰였다. 이 기계는 카드 묶음을 한 번 통과시키며 지정한 한 열에 뚫린 구멍의 숫자에 따라 카드를 칸으로 나누었고, 각 칸 안에서 카드의 순서는 들어온 순서 그대로였으니 한 번의 통과가 곧 한 자리에 대한 안정 정렬이었다. 그래서 여러 자리의 번호로 카드를 정렬하려는 운영자는 가장 낮은 자리부터 한 열씩 통과시키고 칸들을 숫자 순으로 다시 쌓기를 자릿수만큼 되풀이했다. 이 방법과 그 역사, 그리고 컴퓨터 위의 기수 정렬로 이어진 과정은 Knuth의 책에 자세하다\cite{knuth1998taocp3}.
\end{aside}

\subsection{Multiple Keys in stardust's Data Frames}

데이터프레임의 \code{sort\_values(by=[\dots])}는 \cref{alg:multikey}를 그대로 구현한다. 마지막 키부터 한 키씩 안정 정렬하고 단계마다 그 키의 방향(\code{ascending})과 결측값의 위치(\code{na\_position})만 적용한다.

\begin{algorithm}[htbp]
\begin{algorithmic}[1]
\Require 행 순서 $R$, 키 열 $c$, 방향, NA 위치
\State $R$을 $c$의 값이 있는 행 $V$와 NA인 행 $N$으로 가른다 \Comment{각각 순서 유지}
\If{$c$가 문자열}
  \State 고유 문자열을 구하고(해시, 19장) \cref{alg:string-sort}로 정렬해 순위를 매긴다
  \State 행마다 키 $k\gets$ 그 문자열의 순위
\Else
  \State 행마다 키 $k\gets$ \cref{alg:float-key}의 키(정수는 부호 비트 뒤집기, 참·거짓은 한 바이트)
\EndIf
\State \textbf{if} 내림차순 \textbf{then} $k\gets\lnot k$
\State $V$를 $k$로 LSD 기수 정렬 \Comment{안정; 필요한 자리만}
\State \Return NA 위치에 따라 $N\|V$ 또는 $V\|N$
\end{algorithmic}
\caption{데이터프레임 정렬의 한 단계 (stardust \code{sort\_by})}
\label{alg:sort-by}
\end{algorithm}

\cref{alg:sort-by}에서 문자열은 접두 키로 행을 직접 정렬하지 않고 고유 문자열의 순위로 바꾼다. 내림차순이 순위를 뒤집는 것으로 끝나고, 고유값이 적은 열(국가 이름, 범주)에서는 비싼 문자열 정렬을 고유값의 수만큼만 하기 때문이다. 순위는 고유값의 수보다 작으므로 기수 정렬의 자리도 그 수의 바이트만큼만 필요하다. 기수 정렬은 키를 한 번 읽어 할 일을 정한다. 키가 이미 오름차순이면 아무것도 옮기지 않고, 내림차순이면 같은 키의 묶음들을 통째로 거꾸로 놓되 묶음 안의 순서는 지킨다. 그 밖에는 (키$-$최솟값)의 바이트들과 키마다 실제로 다른 바이트들 가운데 통과가 적은 쪽만 정렬하고, 직렬로 많은 키를 정렬할 때는 11비트까지의 넓은 자리를 써서 통과 수를 줄인다. 어느 길이든 결과는 하나의 안정 정렬 순서이므로 데이터는 비용만 정한다. 정렬 결과가 원래 순서와 같으면 행을 옮기는 복사 자체를 건너뛴다. \code{sort\_index}도 인덱스 열을 키로 같은 단계를 쓴다.

\section{Tie-Breaking Rules and Determinism}
\label{sec:ties}

\subsection{Tie-Breaking Rules}

\cref{sec:notation}에서 $\operatorname{arg\,min}$이 집합이라고 했다. 그런데 프로그램은 대개 원소 하나를 돌려주어야 하고, 이때 집합에서 무엇을 고를지 정하지 않으면 답은 구현에 따라 달라진다. 앞에서부터 훑는 반복문은 처음 만난 최솟값을, 뒤에서부터 훑는 반복문은 마지막 최솟값을 돌려주며, 배열을 여러 조각으로 나누어 병렬로 훑는 구현은 조각들의 결과를 어떤 순서로 합치느냐에 따라 또 다른 원소를 돌려줄 수 있다. 같은 일이 정렬에서는 동점의 순서로, 최근접 이웃 찾기에서는 거리가 같은 이웃의 선택으로, 그래프 알고리즘에서는 무게가 같은 변의 선택으로 나타난다.

이 문제의 일관된 해법은 \cref{thm:stable}에서 이미 보았다. 원래의 순서가 약순서라면 거기에 색인처럼 원소마다 다른 값을 키로 덧붙여 사전식 전순서를 만들고 답을 그 전순서로 정의하면 된다. 이렇게 덧붙인 키와 그 순서를 \term{동점 규칙}{tie-breaking rule}이라 부르자. 동점 규칙을 정하면 ``첫 번째 최솟값''처럼 모호했던 말이 ``(값, 색인)의 사전식 순서에서 가장 작은 원소''라는 정의가 되고, 정의를 만족하는 원소는 하나뿐이므로 어떤 알고리즘으로 계산하든 답이 같아야 한다. 이렇게 하면 답은 알고리즘의 실행 경로와 상관없이 \textbf{정의}로 정해진다.

이 원칙은 stardust의 모든 영역에 걸쳐 쓰인다. \code{argmin}과 \code{argmax}는 동점이면 가장 작은 색인을 돌려준다. KD 트리의 이웃 찾기는 (거리, 색인)의 순서로 이웃을 고르고, 트리를 만들 때 점을 나누는 기준도 (좌표, 색인)의 전순서이다. 최소 신장 트리는 변을 (무게, 작은 끝점, 큰 끝점)의 사전식 순서로 비교하며 위상 정렬은 가능한 위상 순서 가운데 사전식으로 가장 작은 것을 돌려준다. 각각의 구현은 그 영역의 장에서 다룬다.

셋째 예는 그래프 이론의 고전적인 사실과 맞닿아 있다. 연결된 그래프의 변 무게가 모두 다르면 최소 신장 트리는 유일하다\cite{kruskal1956,cormen2022clrs}. 증명은 짧다. 서로 다른 최소 신장 트리 $T_1$, $T_2$가 있다고 하고 두 트리의 대칭차에서 무게가 가장 작은 변을 $e$라 하자. $e\in T_1\setminus T_2$라 해도 일반성을 잃지 않는다. $T_2$에 $e$를 더하면 순환이 하나 생기고, 그 순환에는 $T_1$에 없는 변 $f$가 있다($T_1$은 순환을 품지 않으므로). $f$는 대칭차에 속하고 $e$와 다르므로 무게가 모두 다르다는 가정에서 $w(f)>w(e)$이다. 그러면 $T_2+e-f$는 신장 트리이고 $T_2$보다 가벼워 모순이다. 무게가 같은 변이 있으면 이 논증의 마지막 부등식이 깨진다. 그러나 무게에 끝점을 덧붙인 사전식 순서로 변을 비교하면 서로 다른 두 변은 늘 비교되므로 같은 논증이 그 순서에서 성립하고, 답이 하나로 정해진다.

\subsection{Associative Selection and Parallel Computation}

동점 규칙에는 이점이 하나 더 있다. 전순서 위의 최솟값은 계산을 아무렇게나 쪼개고 합쳐도 같은 답을 주므로 병렬 계산이 가능해진다.

\begin{proposition}
\label{prop:min-semilattice}
$(X,<)$가 엄격한 전순서집합이고 $a\wedge b$를 $a$와 $b$ 가운데 작은 쪽이라 하면, 연산 $\wedge$는 결합법칙 $(a\wedge b)\wedge c=a\wedge(b\wedge c)$, 교환법칙 $a\wedge b=b\wedge a$, 멱등법칙 $a\wedge a=a$를 만족한다. 따라서 유한 수열 $x_1,\dots,x_n$에 $\wedge$를 어떤 순서와 어떤 괄호로 적용하든 결과는 $\{x_1,\dots,x_n\}$의 최솟값 하나이다.
\end{proposition}

\begin{proof}
전순서에서는 어떤 두 원소도 비교되므로 $a\wedge b$가 잘 정의되고, 정의에서 교환법칙과 멱등법칙이 바로 나온다. 결합법칙을 보이려면 양변이 모두 $\{a,b,c\}$의 최솟값임을 보이면 된다. $m$을 그 최솟값이라 하자. $(a\wedge b)\wedge c$에서 $m\in\{a,b\}$이면 $a\wedge b=m$이고 $m\le c$이므로 결과는 $m$이며, $m=c$이면 $a\wedge b$가 $c$보다 작을 수 없으므로 역시 결과는 $m$이다. 오른쪽도 같은 논증으로 $m$이다. 마지막 주장은 원소 수에 대한 귀납법으로 따라 나온다. 어떤 괄호를 쳐도 맨 바깥의 연산은 두 부분의 결과를 합치는데, 귀납 가정에서 두 부분의 결과는 각 부분집합의 최솟값이고, 두 최솟값 가운데 작은 쪽은 전체의 최솟값이다.
\end{proof}

\cref{prop:min-semilattice}의 가정에서 전순서가 하는 역할을 분명히 해 두자. 원래의 순서가 약순서일 뿐이면 동점인 두 원소 가운데 ``작은 쪽''이 정해지지 않으므로 구현은 둘 가운데 하나를, 예컨대 왼쪽 피연산자를 고르게 된다. 이 연산은 결합적이지만 교환적이지 않다. 그래서 조각들의 결과를 색인 순서대로 합칠 때에만 직렬 답과 같고, 먼저 끝난 스레드부터 합치는 식으로 피연산자의 순서가 바뀌면 답이 바뀐다. 조각 안과 조각 사이에서 동점일 때 고르는 쪽이 다르면 조각 수에 따라서도 답이 바뀐다(아래 예). 색인을 덧붙여 전순서를 만들면 $\wedge$가 교환적이 되어 이 문제가 사라진다. 조각을 몇 개로 나누든 어느 스레드가 먼저 끝나든 답은 하나이다.

덧셈은 사정이 다르다. 실수의 덧셈은 결합적이고 교환적이지만 부동소수점 수의 덧셈은 반올림 때문에 결합적이지 않아서 합을 나누어 계산하는 방식에 따라 마지막 자리가 달라진다. 그래서 병렬로 계산한 합이 늘 같은 비트를 주게 하려면 덧셈의 순서 자체를 데이터의 길이만으로 고정해야 하는데 그 자세한 내용은 부동소수점 산술을 다루는 장에서 본다. 최솟값과 순위, 정렬처럼 순서만을 쓰는 계산은 전순서만 갖추면 그런 고정 없이도 결과가 하나로 정해지므로 훨씬 다루기 쉽다.

\subsection{\texorpdfstring{\code{argmin} and \code{argmax}}{argmin and argmax} in stardust}

stardust의 \code{argmax}는 ``최댓값 가운데 가장 작은 색인''을 돌려준다. NaN이 있으면 첫 NaN의 색인을 돌려준다. \cref{prop:nan-last}의 순서는 NaN을 $+\infty$보다도 큰 값으로 보므로 여기에 \cref{sec:ties}의 동점 규칙(색인이 작은 쪽)을 더한 전순서에서는 첫 NaN이 최댓값이다. \code{argmin}도 첫 NaN을 돌려준다. NaN이 한 번 나오면 답이 정해진다는 numpy의 관례를 따른 것으로, \code{argmin}에서는 NaN을 가장 작은 값으로 보는 셈이다. \code{nanargmax}는 NaN을 $-\infty$로, \code{nanargmin}은 $+\infty$로 바꿔 읽는다(numpy와 같다).

긴 레인 하나는 여러 스레드에 나누어 계산하는데 결과가 직렬 계산과 같아야 하므로 조각을 합치는 연산을 \cref{prop:min-semilattice}의 전순서 선택으로 만든다.

\begin{algorithm}[htbp]
\begin{algorithmic}[1]
\Require 레인 $x_0,\dots,x_{n-1}$, 조각 길이 $c$
\Ensure \code{argmax}: 첫 NaN의 색인, NaN이 없으면 최댓값의 첫 색인
\State $P\gets\lceil n/c\rceil$
\ForAll{조각 $t=0,\dots,P-1$ \textbf{in parallel}}
  \State $\mathit{at}_t\gets tc+$ (조각 $[tc,\min(n,(t+1)c))$ 안의 직렬 답)
\EndFor
\State $b\gets\mathit{at}_0$
\For{$t=0,\dots,P-1$} \Comment{조각 순서대로}
  \State \textbf{if} $x_{\mathit{at}_t}$가 NaN \textbf{then return} $\mathit{at}_t$ \Comment{앞 조각의 NaN이 이긴다}
  \State \textbf{if} $t>0$ \textbf{and} $x_{\mathit{at}_t}>x_b$ \textbf{then} $b\gets\mathit{at}_t$ \Comment{같으면 앞 조각}
\EndFor
\State \Return $b$
\end{algorithmic}
\caption{조각으로 나눈 결정적인 argmax (stardust \code{arg\_whole})}
\label{alg:argmax-parallel}
\end{algorithm}

\cref{alg:argmax-parallel}이 직렬 답과 같다는 것을 보이자. 조각 안의 답 $\mathit{at}_t$는 그 조각에서 전순서 ``NaN이 가장 크고, 값이 클수록 크고, 같으면 색인이 작을수록 크다''의 최댓값이다. 합치는 루프는 같은 전순서로 조각 답들의 최댓값을 고른다. NaN을 만나면 그 조각보다 앞에는 NaN이 없으므로(앞 조각의 답이 NaN이 아니었으므로) 그것이 첫 NaN이고, NaN이 없으면 값이 엄격히 클 때만 바꾸므로 같은 값이면 앞 조각, 곧 작은 색인이 남는다. 전순서 위의 최댓값은 \cref{prop:min-semilattice}에 따라 어떻게 묶어도 같으므로 결과는 레인 전체의 최댓값이고 조각의 수와 스레드 수에 의존하지 않는다.

동점 규칙을 빠뜨리면 어떻게 되는지 직접 보자. 값 $1,2,3,5$가 섞인 길이 $40$의 배열에 $5$가 21번 나온다. 조각 안에서는 첫 최댓값을 고르되 조각을 합칠 때 값이 같으면 오른쪽 조각의 답을 고르는(동점 규칙이 없는) 구현은 조각 수가 $1,2,3,4,5,8$일 때 답이 $1,21,29,30,32,35$로 모두 다르다(\code{figures/01/argmax\_chunks.py}). (값, 색인)의 전순서로 합치면 언제나 $1$이다(\cref{fig:argmax-tree}).

\begin{figure}[htbp]
\centering
\begin{tikzpicture}[x=1.05cm,y=0.9cm,font=\small,every node/.style={inner sep=2pt}]
\foreach \i/\lab in {0/{$(5,1)$},1/{$(5,12)$},2/{$(5,21)$},3/{$(5,30)$}}{\node[draw,rounded corners=1pt] (l\i) at (\i*2.2,0) {\lab};}
\node[draw,rounded corners=1pt] (m0) at (1.1,1.4) {$(5,1)$};
\node[draw,rounded corners=1pt] (m1) at (5.5,1.4) {$(5,21)$};
\node[draw,rounded corners=1pt,fill=gray!20] (r) at (3.3,2.8) {$(5,1)$};
\draw (l0) -- (m0) -- (l1); \draw (l2) -- (m1) -- (l3); \draw (m0) -- (r) -- (m1);
\node[anchor=west,align=left] at (7.6,1.4) {조각마다 (값, 색인)의 답\\합칠 때: 값이 크면 이기고\\같으면 색인이 작은 쪽};
\end{tikzpicture}
\caption{(값, 색인)의 전순서로 조각의 답을 합치는 argmax. 어떤 순서와 괄호로 합쳐도 루트의 답은 $(5,1)$이다.}
\label{fig:argmax-tree}
\end{figure}


조각 안의 직렬 답은 \code{float64}의 연속 레인에서 데이터를 두 번 읽어 구한다. 첫 번째 읽기는 극값만 구한다. 레지스터 네 개에 독립된 \code{FMAX} 사슬 넷을 두고 원소 여덟 개씩 접는다. \code{FMAX}는 NaN을 전파하고 $-0<+0$으로 다루는 명령이다. 사슬이 넷인 것은 명령 하나의 지연 시간 동안 다른 사슬의 명령을 돌려 산술 장치를 쉬지 않게 하려는 것이다. 두 번째 읽기는 그 값과 \code{==}로 같은 첫 색인을 찾는다. 극값이 NaN이면 첫 NaN을 찾는다. 위치를 함께 따라가며 한 번에 읽는 방법은 원소마다 앞의 선택에 의존하는 사슬이 되어 더 느렸다.

\begin{algorithm}[htbp]
\begin{algorithmic}[1]
\Require 연속 레인 $x_0,\dots,x_{n-1}$ ($n\ge1$)
\Ensure \code{argmax}의 직렬 답
\State $m_0,\dots,m_3\gets$ 처음 여덟 원소(레지스터마다 둘) \Comment{$n<8$이면 스칼라로}
\For{$i=8,16,\dots$}
  \State $m_r\gets\mathrm{FMAX}(m_r,x_{i+2r..i+2r+1})$ ($r=0,\dots,3$) \Comment{독립된 사슬 넷}
\EndFor
\State $b\gets$ 네 레지스터와 남은 원소의 $\mathrm{FMAX}$
\If{$b$가 NaN} \Return 첫 NaN의 색인 \EndIf
\State \Return $x_k=b$인 첫 $k$ \Comment{넷씩 비교해 찾으면 멈춘다}
\end{algorithmic}
\caption{두 번 읽는 argmax 커널 (stardust \code{neon\_arg\_extreme})}
\label{alg:argmax-kernel}
\end{algorithm}

\cref{alg:argmax-kernel}이 ``최댓값의 첫 색인''을 내는지는 $\pm0$에서 따져 볼 만하다. \code{FMAX}는 $-0$과 $+0$ 가운데 $+0$을 고르지만 두 번째 읽기의 \code{==}는 $-0$과 $+0$을 같게 보므로 처음 나온 영의 색인을 돌려준다. 최댓값이 영이고 $-0$이 먼저 나오는 레인 $(-0,+0)$에서 답은 $0$이다. 직렬로 엄격한 비교 $x_i>x_b$를 써서 첫 최댓값을 따라가도 두 영은 서로 크지 않으므로 답은 $0$이다. 두 방법은 같은 계약을 낸다. \code{nanargmax}는 \code{FMAXNM}(NaN을 무시하는 명령)을 쓰고 NaN을 $-\infty$로 읽으므로 실제 $-\infty$와 NaN이 동점이 되어 색인이 작은 쪽이 이긴다. 모든 원소가 NaN이면 고를 것이 없으므로 \code{ValueError}이다.

축이 마지막 축이 아닌 경우(\code{argmax(m, axis=0)})에는 레인이 메모리를 가로질러 놓인다. 레인마다 띄엄띄엄 읽는 대신 행을 차례로 읽으며 열 $256$개씩의 현재 최댓값과 색인을 함께 갱신한다. 엄격한 비교로만 바꾸므로 첫 색인이 남고, NaN은 그 열에 아직 NaN이 없을 때만 자리를 차지하므로 첫 NaN이 남는다. 레인 커널과 같은 답이다. 길이가 $2^{15}$ 이상인 레인 하나는 스레드 풀이 있으면 길이 $2^{12}$의 조각으로 나누어 \cref{alg:argmax-parallel}을 쓴다.

\section{Equivalence in Practice: \texorpdfstring{\code{unique}}{unique}}
\label{sec:unique}

\cref{thm:partition}을 구현하는 대표 함수가 \code{unique}이다. stardust에는 \code{unique\_values}, \code{unique\_counts}, \code{unique\_inverse}, \code{unique\_all}의 넷이 있고 모두 같은 동치관계, 곧 모든 NaN을 한 값으로 보고 $-0$과 $+0$도 한 값으로 보는 동치관계를 쓴다. 동치류를 대표해 보여 주는 값은 그 동치류에서 입력 순서로 처음 나온 원소이다. 그래서 $-0$이 먼저 나왔으면 대표는 $-0$이고 NaN의 대표는 첫 NaN의 비트 그대로이다. 결과는 정렬된 순서이고 NaN의 동치류는 끝에 온다.

\begin{algorithm}[htbp]
\begin{algorithmic}[1]
\Require 배열 $x$ (C 순서로 편 길이 $n$)
\Ensure 대표값 $v$, 첫 색인 $\mathit{idx}$, 개수 $\mathit{cnt}$, 역사상 $\mathit{inv}$
\State 쌍 $(x_i,i)$를 \cref{alg:split-nans}로 갈라 NaN이 아닌 $m$개를 앞에
\State 앞의 $m$개를 안정 정렬 \Comment{(값, 색인)의 전순서}
\State $S\gets\{0\}$ \Comment{동치류가 시작하는 정렬 위치}
\For{$i=1,\dots,n-1$}
  \State \textbf{if} ($i<m$ \textbf{and} 값$_{i-1}\ne$ 값$_i$) \textbf{or} $i=m$ \textbf{then} $S\gets S\cup\{i\}$ \Comment{$\ne$는 $\pm0$을 같게 본다}
\EndFor
\ForAll{동치류 $k$: 정렬 위치 $[s_k,s_{k+1})$ ($s_K=n$; $n=0$이면 빈 결과)}
  \State $v_k\gets$ 값$_{s_k}$; $\mathit{idx}_k\gets$ 색인$_{s_k}$; $\mathit{cnt}_k\gets s_{k+1}-s_k$
  \State 위치 $i\in[s_k,s_{k+1})$마다 $\mathit{inv}[\text{색인}_i]\gets k$
\EndFor
\end{algorithmic}
\caption{정렬로 구하는 unique (stardust \code{unique\_parts})}
\label{alg:unique}
\end{algorithm}

\cref{alg:unique}의 정확성은 세 결과를 합친 것이다. \cref{prop:sorted-unique}에 따라 동치류는 정렬 순서에서 연속한 구간이므로 이웃끼리만 비교하면 구간의 경계를 찾는다. 안정 정렬이므로 구간의 첫 원소는 그 동치류에서 입력 순서로 처음 나온 원소이고(\cref{thm:stable}), 그것이 대표와 첫 색인이 된다. NaN들은 1행에서 입력 순서대로 끝에 모였으므로 하나의 구간이 되고 그 첫 원소가 첫 NaN이다. 5행은 비교 $\ne$를 NaN이 아닌 값들에만 쓴다. NaN에 \code{!=}를 쓰면 모든 NaN이 저마다 새 동치류가 되기 때문이다. 마지막 루프는 동치류 단위로 나누어 병렬로 돌 수 있다. 원소마다 정확히 한 동치류에 속하므로 $\mathit{inv}$의 각 칸은 한 번만 쓰인다.

\code{unique\_values}처럼 색인이 필요 없으면 쌍 대신 값만 안정 정렬한다. 옮길 데이터가 절반이고, \cref{alg:zero-signs} 덕분에 값만 정렬해도 $\pm0$의 대표가 입력 순서로 처음 나온 영이 된다. 정렬된 값을 앞에서부터 훑으며 구간의 첫 값만 남기면 끝이다. 시간은 정렬이 지배하므로 $O(n\log n)$(기수 정렬이면 $O(n)$)이다. 데이터프레임의 \code{unique}와 \code{groupby}는 결과를 입력 순서(처음 나온 순서)로 돌려주므로 정렬 대신 해시를 쓰며, 그 구현은 19장과 42장에서 다룬다.


\section{Ranks}
\label{sec:ranks}

통계에서는 정렬 그 자체보다 각 관측값이 몇 번째인지를 나타내는 순위가 필요한 경우가 많다. 순위는 동점을 어떻게 다루느냐에 따라 여러 가지로 정의되며 순위를 이용하는 검정 통계량은 이 선택에 따라 값이 달라진다. $\pi$를 안정 정렬 순열이라 하고, 원소 $x_i$가 정렬된 수열에서 놓인 위치를 $p_i=\pi^{-1}(i)$라 하자. 또 $x_i$의 동치류가 \cref{prop:sorted-unique}에서처럼 위치 $c_i+1,\dots,c_i+e_i$를 차지한다고 하자.

\begin{definition}[순위]
\label{def:ranks}
$x_i$의 순위를 다음과 같이 정의한다.
\begin{itemize}
\item \term{서수 순위}{ordinal rank}: $p_i$.
\item \term{최소 순위}{min rank}: $c_i+1$.
\item \term{최대 순위}{max rank}: $c_i+e_i$.
\item \term{평균 순위}{average rank}: $c_i+(e_i+1)/2$.
\item \term{조밀 순위}{dense rank}: $x_i$보다 엄격히 작은 값들이 이루는 동치류의 수에 $1$을 더한 값.
\end{itemize}
\end{definition}

동점이 없으면 다섯 순위가 모두 같다. 평균 순위는 동점인 원소들이 차지한 위치 $c_i+1,\dots,c_i+e_i$의 평균으로, 통계학에서 중간 순위(mid-rank)라고도 부르며 Kendall 이래 순위 통계에서 동점을 다루는 표준적인 방법이다\cite{kendall1948rank}. 최소 순위는 운동 경기의 순위와 같아서 공동 2위가 둘이면 다음 선수는 4위가 되고 조밀 순위에서는 3위가 된다. \cref{tab:ranks}는 일곱 개의 관측값에 대해 다섯 순위를 계산한 것이다. 이 다섯 가지는 SciPy의 \code{rankdata}가 제공하는 방법과 같다\cite{scipy-rankdata}.

\begin{table}[htbp]
\begin{sdtable}{colspec={cccccccc}, row{1}={font=\small\bfseries}}
\toprule
$i$ & 1 & 2 & 3 & 4 & 5 & 6 & 7 \\
\midrule
$x_i$ & 3.2 & 1.5 & 3.2 & 0.7 & 3.2 & 1.5 & 2.8 \\
서수 & 5 & 2 & 6 & 1 & 7 & 3 & 4 \\
최소 & 5 & 2 & 5 & 1 & 5 & 2 & 4 \\
최대 & 7 & 3 & 7 & 1 & 7 & 3 & 4 \\
평균 & 6 & 2.5 & 6 & 1 & 6 & 2.5 & 4 \\
조밀 & 4 & 2 & 4 & 1 & 4 & 2 & 3 \\
\bottomrule
\end{sdtable}
\caption{동점이 있는 일곱 관측값의 다섯 가지 순위. $3.2$는 위치 $5$, $6$, $7$을, $1.5$는 위치 $2$, $3$을 차지한다.}
\label{tab:ranks}
\end{table}

평균 순위가 표준이 된 데에는 이유가 있다. 순위의 합은 많은 순위 검정 통계량의 재료이고 평균 순위를 쓰면 그 합이 동점의 유무와 상관없이 같은 값으로 유지된다.

\begin{proposition}
\label{prop:rank-sum}
평균 순위를 $r_i$라 하면 $\sum_{i=1}^{n}r_i=n(n+1)/2$이다. 또한 서수 순위의 합도 같은 값이고, 최소 순위의 합은 이보다 작거나 같으며 최대 순위의 합은 이보다 크거나 같다. 두 부등식에서 등호는 동점이 없을 때에만 성립한다.
\end{proposition}

\begin{proof}
서수 순위 $p_1,\dots,p_n$은 $1,\dots,n$을 한 번씩 늘어놓은 것이므로 합이 $n(n+1)/2$이다. 동치류 하나가 위치 $c+1,\dots,c+e$를 차지한다고 하면, 그 원소들의 서수 순위의 합은 $\sum_{t=1}^{e}(c+t)=ec+e(e+1)/2$이고, 평균 순위의 합은 $e\,(c+(e+1)/2)$로 이와 같다. 동치류별로 합이 같으므로 전체 합도 같다. 같은 동치류에서 최소 순위의 합은 $e(c+1)$이고 최대 순위의 합은 $e(c+e)$이며, 이 둘과 서수 순위의 합의 차는 각각 $e(e-1)/2$이다. 이 값은 $e\ge2$일 때 양수이고 $e=1$일 때 $0$이므로 나머지 주장이 따라 나온다.
\end{proof}

순위는 안정 정렬 한 번과 선형 시간의 훑기로 계산할 수 있다. 정렬된 순서로 원소를 훑으며 동점인 원소가 이어지는 구간을 찾고 구간마다 그 위치들로 순위를 매기면 된다(\cref{alg:ranks}).

\begin{algorithm}[htbp]
\begin{algorithmic}[1]
\Require 수열 $x_0,\dots,x_{n-1}$과 엄격한 약순서 $<$
\Ensure 평균 순위 $r_0,\dots,r_{n-1}$
\State $\pi\gets$ $x$의 안정 정렬 순열 \Comment{$x_{\pi(0)},\dots,x_{\pi(n-1)}$이 정렬된 수열}
\State $s\gets0$
\While{$s<n$}
  \State $t\gets s+1$
  \While{$t<n$ \textbf{and} $x_{\pi(t)}\sim x_{\pi(s)}$} \Comment{동점 구간의 끝을 찾는다}
    \State $t\gets t+1$
  \EndWhile
  \For{$q\gets s$ \textbf{to} $t-1$} \Comment{위치 $s+1,\dots,t$의 평균}
    \State $r_{\pi(q)}\gets(s+1+t)/2$
  \EndFor
  \State $s\gets t$
\EndWhile
\State \Return $r$
\end{algorithmic}
\caption{안정 정렬을 이용한 평균 순위의 계산}
\label{alg:ranks}
\end{algorithm}

\cref{alg:ranks}는 각 원소를 구간의 첫 원소하고만 견준다. \cref{prop:sorted-unique}에 따라 동치류는 연속한 구간을 이루고 $\sim$이 추이적이므로 이것으로 구간의 끝을 정확히 찾는다(이웃끼리 견주어도 같다). 다만 이 비교 역시 약순서의 동점이어야 한다. NaN이 섞인 부동소수점 수에서 \code{==}를 그대로 쓰면 NaN들이 저마다 다른 구간이 되어 버린다.

\subsection{Ranks in stardust}

\code{sd.stats.rankdata}는 \cref{def:ranks}의 다섯 방법(\code{average}, \code{min}, \code{max}, \code{dense}, \code{ordinal})을 모두 제공하고 레인마다 \cref{alg:ranks}를 돈다. 안정 정렬 대신 (값, 위치)의 쌍을 정렬하는데, 쌍의 사전식 비교가 바로 \cref{thm:stable}의 $\prec$이므로 어떤 정렬 알고리즘을 써도 결과는 안정 정렬과 같다. 동점 구간 $[s,t)$를 찾으면 그 안의 원소마다 방법에 따라 순위를 쓴다. \code{average}는 $(s+1+t)/2$, \code{min}은 $s+1$, \code{max}는 $t$, \code{dense}는 지금까지 지난 구간의 수, \code{ordinal}은 정렬 위치에 $1$을 더한 값이다. 결과는 모두 \code{float64}이다. 평균 순위 $(s+1+t)/2$는 정수의 합을 $2$로 나눈 것이어서 $n<2^{53}$이면 반올림 없이 정확하다. 합이 $2^{53}$ 아래의 정수이고 $2$로 나누는 것은 지수만 바꾸기 때문이다.

NaN의 처리는 \code{nan\_policy}가 정한다. 기본값 \code{'propagate'}이면 NaN이 하나라도 있는 레인은 순위 전체가 NaN이고, \code{'omit'}이면 NaN의 자리만 NaN으로 두고 나머지를 그들끼리 순위를 매기며, \code{'raise'}이면 \code{ValueError}이다. NaN을 어디에 둘지 순서를 정하는 대신 순위에서 빼는 것이다. 순위의 합이 $n(n+1)/2$라는 \cref{prop:rank-sum}을 쓰는 검정 통계량에서 NaN에 순위를 주면 그 합이 뜻을 잃기 때문이다.

NaN이 아닌 원소가 $2^{16}$개 이상인 레인은 표본 정렬(samplesort)로 여러 스레드에 나누어 정렬한다\cite{frazer1970samplesort}.

\begin{algorithm}[htbp]
\begin{algorithmic}[1]
\Require 쌍의 배열 $v[0..n)$ ($n\ge2^{16}$), 통의 수 $64$, 과표본 $16$, 조각 길이 $2^{14}$
\State $v$에서 고르게 떨어진 $64\cdot16$개를 표본으로 뽑아 정렬
\State 분할값 $s_b\gets$ 표본의 $16b$번째 ($b=1,\dots,63$)
\ForAll{조각 $p$ \textbf{in parallel}}
  \For{조각의 원소 $x$}
    \State $\mathit{id}(x)\gets|\{b:s_b<x\}|$ \Comment{가지 없는 이분 탐색}
    \State $c[p][\mathit{id}(x)]\gets c[p][\mathit{id}(x)]+1$
  \EndFor
\EndFor
\State (통, 조각)의 순서로 $c$를 누적해 조각마다 통별 첫 자리를 정한다
\ForAll{조각 $p$ \textbf{in parallel}} 원소를 자기 통의 자리에 순서대로 흩뿌린다 \EndFor
\ForAll{통 $b$ \textbf{in parallel}} 통 안을 정렬 \EndFor
\end{algorithmic}
\caption{표본 정렬 (stardust \code{sort\_long}, rankdata의 긴 레인)}
\label{alg:samplesort}
\end{algorithm}

\cref{alg:samplesort}의 결과가 스레드 수와 무관한 까닭은 결정에 쓰이는 모든 것이 데이터와 길이로만 정해지기 때문이다. 표본의 위치, 분할값, 조각의 경계, 통 안의 순서(조각 순서와 조각 안의 순서)가 모두 그렇다. 게다가 쌍의 순서는 전순서이므로 정렬된 결과 자체가 하나뿐이다. 균형도 이 전순서에서 나온다. 모든 값이 같은 데이터라도 쌍은 위치로 갈리므로 분할값이 위치로 고르게 나뉘어, 한 통에 원소가 몰리지 않는다. 과표본을 $16$배로 잡는 것은 표본의 분위수가 실제 분위수에서 크게 벗어날 확률을 줄이려는 것으로, 통의 크기는 높은 확률로 평균 $n/64$의 상수배 안에 든다.

\section{Strings and Lexicographic Keys}
\label{sec:string-keys}

문자열의 순서는 길이가 다른 수열의 사전식 순서(\cref{sec:relations})이고 stardust는 UTF-8 바이트의 순서로 비교한다. UTF-8은 바이트 순서가 코드 포인트 순서와 같도록 설계되었으므로(15장) 이 순서는 곧 유니코드 코드 포인트의 사전식 순서이다. 문자열을 비교 함수로 정렬하면 비교 한 번이 바이트 여럿을 읽고 문자열마다 다른 메모리를 찾아가야 해서 느리다. 데이터프레임은 대신 문자열을 64비트 정수 키로 바꿔 기수 정렬한다.

\begin{proposition}
\label{prop:string-prefix}
문자열 $s$의 앞 8바이트를 첫 바이트가 가장 높은 자리가 되도록(빅 엔디언) 읽고 8바이트가 안 되면 남는 자리를 $0$으로 채운 부호 없는 정수를 $p(s)$라 하자. 두 문자열 $s,t$에 대해 $p(s)<p(t)$이면 $s$는 사전식으로 $t$보다 앞선다.
\end{proposition}

\begin{proof}
$p(s)<p(t)$이면 두 수의 바이트를 높은 자리부터 비교할 때 처음 다른 자리 $j<8$에서 $s$의 바이트(또는 채운 $0$)가 작다. $s$의 바이트가 실제 바이트이면 $t$의 $j$번째 바이트도 그보다 크므로 실제 바이트이고, $j$ 앞의 바이트가 같으니 사전식 순서에서 $s$가 앞선다. $s$가 $j$에서 이미 끝났으면(채운 $0$) $s$는 $t$의 앞부분이고 $t$가 더 길므로 역시 $s$가 앞선다.
\end{proof}

역은 성립하지 않는다. $p(s)=p(t)$인 두 문자열은 앞 8바이트가 같거나, 짧은 쪽에 채운 $0$과 실제 바이트 $0$이 같은 자리에 온 경우(\code{"a"}와 \code{"a\textbackslash0"})이다. 그래서 키가 같은 행들의 구간마다 다음 8바이트로 같은 일을 되풀이한다(\cref{fig:string-keys}). 구간이 짧거나 모든 문자열이 끝났으면 바이트를 직접 비교하는 정렬로 마무리하며 이 마지막 비교가 채운 $0$과 실제 $0$을 길이로 가린다.

\begin{algorithm}[htbp]
\begin{algorithmic}[1]
\Function{SortStrings}{행 목록 $R$, 깊이 $d$} \Comment{$R$의 행은 앞 $d$바이트를 $0$으로 채워 견주면 같다}
  \If{$|R|$이 짧음 \textbf{or} $R$의 모든 문자열이 $d$바이트 안에서 끝남}
    \State $R$을 $d$바이트부터의 사전식 비교로 안정 정렬 \Comment{길이까지 본다}
    \State \Return
  \EndIf
  \State 행마다 키 $k\gets$ 바이트 $d..d+7$의 빅 엔디언 정수(모자라면 $0$)
  \State $R$을 $k$로 LSD 기수 정렬 \Comment{안정}
  \ForAll{$k$가 같은 연속 구간 $R'$}
    \State \Call{SortStrings}{$R'$, $d+8$}
  \EndFor
\EndFunction
\end{algorithmic}
\caption{접두 키로 문자열 정렬 (데이터프레임의 문자열 정렬, MSD 방식)}
\label{alg:string-sort}
\end{algorithm}

\begin{figure}[htbp]
\centering
\begin{tikzpicture}[x=1cm,y=0.6cm,font=\small]
\node[font=\small\bfseries] at (0.9,1) {입력};
\node[font=\small\bfseries] at (4.6,1) {앞 8바이트 키로};
\node[font=\small\bfseries] at (8.6,1) {키가 같은 구간을 다음 8바이트로};
\foreach \s [count=\i from 0] in {banana,applesaurus,apple,applesauce,apricot}{\node[anchor=west] at (0,-\i) {\texttt{\s}};}
\foreach \s/\k [count=\i from 0] in {apple/{apple\textperiodcentered\textperiodcentered\textperiodcentered},applesaurus/applesau,applesauce/applesau,apricot/{apricot\textperiodcentered},banana/{banana\textperiodcentered\textperiodcentered}}{\node[anchor=west] at (3.2,-\i) {\texttt{\k}};}
\draw[gray!70] (3.15,-0.55) rectangle (5.0,-2.45);
\foreach \s [count=\i from 0] in {apple,applesauce,applesaurus,apricot,banana}{\node[anchor=west] at (7.4,-\i) {\texttt{\s}};}
\node[anchor=west,font=\scriptsize] at (5.05,-1.5) {\texttt{rus} 대 \texttt{ce}};
\draw[->] (2.4,-2) -- (3.1,-2); \draw[->] (6.2,-2) -- (7.3,-2);
\end{tikzpicture}
\caption{8바이트 접두 키의 정렬. 첫 단계에서 키가 같은 행(상자)만 다음 8바이트로 다시 정렬한다. \texttt{applesaurus}와 \texttt{applesauce}는 앞 8바이트 \texttt{applesau}가 같아서 둘째 단계의 \texttt{rus}와 \texttt{ce}로 갈린다. 가운데 점은 채운 $0$ 바이트이다.}
\label{fig:string-keys}
\end{figure}

\cref{alg:string-sort}는 \cref{thm:lsd}와 반대로 가장 중요한 바이트부터 정렬하는 MSD 방식이다. 첫 8바이트 키로 정렬하면 \cref{prop:string-prefix}에 따라 키가 작은 행은 사전식으로도 앞선다. 키가 같은 행들은 \cref{prop:sorted-unique}에 따라 연속한 구간을 이루고, 그 구간 안의 순서는 나머지 바이트와 길이가 정하므로 구간마다 같은 일을 되풀이하면 \cref{thm:lex}의 (첫 키, 나머지)의 사전식 순서를 얻는다. 각 단계의 기수 정렬과 마지막 비교 정렬이 안정하므로 끝까지 동점인 행은 입력 순서를 지킨다.

비용은 공통 접두사의 길이가 정한다. 키가 같은 구간의 공통 접두사가 $L$바이트이면 그 구간은 깊이 $\lceil L/8\rceil$까지 내려가고, 깊이마다 구간의 원소들이 기수 정렬 한 번(필요한 바이트만큼의 통과)을 겪는다. $128$행보다 짧은 구간은 기수 정렬의 고정비가 아까우므로 바이트 비교 정렬로 끝낸다. 앞 8바이트로 대부분 갈리는 문자열(이름, 코드)은 기수 정렬 한 번이면 끝나고, 공통 접두사가 긴 문자열(경로, URL)은 깊이 들어갈수록 구간이 좁아진다. $2^{17}$행이 넘는 구간은 호출한 스레드가 다음 8바이트를 병렬 기수 정렬로 정렬하고 다시 내려간다. 바이트 비교에서 주의할 점이 하나 있다. C++의 \code{char}는 부호가 있을 수 있어서 \code{char}끼리 비교하면 \code{0x80} 이상의 바이트(UTF-8의 다바이트 문자)가 음수가 되어 ASCII보다 앞서 버린다. stardust는 바이트를 부호 없는 값으로 비교하는 \code{std::string\_view}의 비교(\code{char\_traits<char>::compare})를 쓴다.

UTF-8의 바이트 순서가 코드 포인트 순서와 같은 까닭도 짧게 보일 수 있다. 코드 포인트는 크기에 따라 $1$--$4$바이트로 적히고, 첫 바이트의 범위가 길이에 따라 \code{00}--\code{7F}, \code{C2}--\code{DF}, \code{E0}--\code{EF}, \code{F0}--\code{F4}로 늘어나므로 길이가 다른 두 문자는 첫 바이트에서 코드 포인트의 순서대로 갈린다. 길이가 같으면 코드 포인트의 비트가 고정된 머리 비트 뒤에 높은 자리부터 차례로 채워지므로 바이트열의 사전식 순서가 코드 포인트의 크기 순서이다. 한 문자의 부호화는 다른 문자의 부호화의 앞부분이 되지 않으므로, 두 문자열의 바이트를 비교하다 처음 다른 바이트는 처음 다른 코드 포인트 안에 있고 그 비교가 그 코드 포인트의 순서와 같다. 자세한 부호화 규칙은 15장에서 다룬다.



\section{Testing the Order Contracts}
\label{sec:order-testing}

순서와 동치의 계약은 값이 아니라 위치와 비트로 드러나므로 시험도 그것을 본다. 기준(oracle)은 계약을 그대로 옮긴 느린 구현이다. 예를 들어 \code{argsort}의 기준은 Python의 안정 정렬에 키 (NaN인가, 값)을 준 것(내림차순이면 (NaN이 아닌가, $-$값))이다. 이 키 자체가 \cref{prop:nan-last}의 약순서이고 Python 정렬이 안정하므로 기준은 계약의 정의와 같다. \cref{tab:order-tests}는 시험이 지나는 경계와 그 이유를 모은 것이다.

\begin{table}[htbp]
\begin{sdtable}{colspec={lX}, row{1}={font=\small\bfseries}}
\toprule
시험 & 무엇을 왜 \\
\midrule
경로의 경계 & 길이 $16/17$(네트워크), $31/32/33$(NEON 정렬), $127/128$(쌍의 기수 정렬), $1023/1024$(값의 기수 정렬), $2^{17}$ 근처(조각)의 양쪽을 지나 모든 경로를 시험한다 \\
동점이 많은 값 & 정수 $-5..5$, NaN $10\%$, $-0$ $5\%$를 섞은 표본; 안정성과 $\pm0$의 순서가 드러난다 \\
비트 비교 & 결과를 64비트 정수로 보고 견준다: $-0$과 $+0$, NaN의 비트가 다르면 실패 \\
대표값 & \code{unique}의 대표가 동치류의 첫 원소인지($\pm0$이 섞인 경우, 비트가 다른 NaN들) \\
스레드 수 & 배열은 $1,2,3,7,10$, \code{rankdata}는 $1,3,10$에서 같은 바이트인지 \\
numpy와의 차이 & 같은 계약은 결과를 견주고, 다른 곳(예: numpy의 \code{unique}는 \code{[0.0, -0.0]}에서 $-0$을 대표로 돌려준다, \code{figures/01/numpy\_unique.py})은 차이를 문서에 적고 우리 계약대로 확인한다 \\
\bottomrule
\end{sdtable}
\caption{순서 계약의 시험.}
\label{tab:order-tests}
\end{table}

\section{APIs in This Chapter}
\label{sec:order-apis}

이 장은 특정 영역의 API보다 여러 영역이 함께 지키는 계약을 다루었다. 정렬의 안정성과 NaN의 위치, 동치의 정의와 대표값, 동점 규칙이 그것이다. \cref{tab:order-apis}은 그 계약을 직접 드러내는 API와 매개변수, 그리고 구현을 설명한 곳을 모은 것이다. 같은 순서 규칙을 쓰는 API가 더 있다. \code{searchsorted}는 정렬된 배열에서 값의 동치류가 차지하는 구간의 양 끝, 곧 \cref{prop:sorted-unique}의 $c$와 $c+e$를 \code{side='left'}와 \code{'right'}로 돌려준다(31장). \code{stats.mode}는 가장 잦은 값이 여럿이면 가장 작은 값을 고른다(111장). 정렬 알고리즘 자체는 18장과 31장에서, 해시를 쓰는 묶기는 19장과 42장에서 다룬다.

\begin{table}[htbp]
\begin{sdtable}{colspec={lXl}, row{1}={font=\small\bfseries}}
\toprule
API & 계약과 매개변수 & 설명한 곳 \\
\midrule
\code{sort}, \code{argsort} & 안정(\code{stable=True}), NaN 끝(\code{descending=True}면 앞), $\pm0$의 입력 순서 & \cref{sec:comparators,sec:float-order,sec:sorting,sec:long-lane} \\
\code{argmin}, \code{argmax} & 첫 NaN, 동점이면 작은 색인, 스레드와 무관 & \cref{sec:ties}, \cref{alg:argmax-parallel,alg:argmax-kernel} \\
\code{nanargmin}, \code{nanargmax} & NaN을 $\pm\infty$로, 모두 NaN이면 \code{ValueError} & \cref{sec:ties} \\
\code{unique\_*} & NaN 하나, $\pm0$ 하나, 대표는 첫 원소, 정렬된 순서 & \cref{sec:unique}, \cref{alg:unique} \\
\code{stats.rankdata} & \code{method} 다섯 가지, \code{nan\_policy} 셋 & \cref{sec:ranks}, \cref{alg:ranks,alg:samplesort} \\
\code{DataFrame.sort\_values}, \code{sort\_index} & 여러 키, 키별 \code{ascending}, \code{na\_position}, 문자열은 UTF-8 바이트 순서 & \cref{sec:sorting,sec:string-keys}, \cref{alg:multikey,alg:sort-by,alg:string-sort} \\
\bottomrule
\end{sdtable}
\caption{이 장의 계약을 직접 드러내는 stardust API.}
\label{tab:order-apis}
\end{table}

\section{Concluding Remarks}

이 장의 내용을 이론, 알고리즘, 구현으로 나누어 돌아보자. 이론에는 비교 가능한 것과 같은 것을 구별하는 약순서와 동치관계가 있고 모든 약순서는 키로 비교하는 순서라는 표현 정리가 있다. 알고리즘에서는 그 약순서에 색인 같은 키를 덧붙여 전순서를 만드는 것이 거의 모든 문제를 푼다. 안정 정렬, 여러 키의 정렬, 결정적인 동점 규칙, 병렬 축약의 결정성이 모두 이 생각에서 나온다. 구현에서는 부동소수점 수가 문제를 일으킨다. NaN은 약순서를 깨고 $\pm0$은 동점이면서 구별되는 값이어서 NaN을 미리 떼어 내고 영의 부호를 기록하고 비트를 정수 키로 바꾸는 손질이 필요하다. 다음 장에서는 그 정수 키의 바탕인 정수와 나머지 산술, 2의 보수를 다룬다.

\chapterreferences
